\documentclass[10pt]{article}
\usepackage[a4paper,margin=2cm,landscape]{geometry}
\usepackage[T1]{fontenc}
\usepackage{lmodern}
\input{glyphtounicode}\pdfgentounicode=1

\usepackage{booktabs}
\usepackage{graphicx}
\usepackage{longtable}
\usepackage{float}
\usepackage{natbib}
\usepackage{amsmath}
\usepackage{amssymb}
\newcommand{\Cs}{C}
\renewcommand{\thetable}{S\arabic{table}}
\title{Online Resource 1\\[0.5ex]\large Characterizing Scanpath Decoding Pipelines with Perturbation Profiles\\[0.5ex]\normalsize Saleh Mozaffari and Abolfazl Nadi}
\author{}
\date{}
\begin{document}
\maketitle
\noindent Tables S1--S13: every profile of the primary analysis (within-fold AUC, class-balanced participant folds); $R = 6$ replicates, $B = 60$ refit-bootstrap draws, seed 42; linear SVM, ten-fold participant-held-out cross-validation; duration-marginal donor pool of the other participants' durations. $u$ is the reference scale (one text line in SB-SAT and EMIP, $2^\circ$ in GazeBase). The same profiles with the pooled AUC are in the released outputs.
\begin{table}[!htbp]\centering\scriptsize
\caption{Profile of \texttt{gb\_tex\_vd1}: reading vs. video viewing. Sensitivity (AUC drop) with refit-bootstrap 95\% interval; bold: interval excludes zero.}
\resizebox{\linewidth}{!}{\begin{tabular}{lrllllllll}
\toprule
Descriptor & AUC & order & dur.\ coupl. & dur.\ marg. & rotation & jit.\ 0.25$u$ & jit.\ 0.5$u$ & jit.\ 1$u$ & jit.\ 2$u$ \\
\midrule
Count & 0.969 & 0 (ex.) & 0 (ex.) & 0 (ex.) & 0 (ex.) & 0 (ex.) & 0 (ex.) & 0 (ex.) & 0 (ex.) \\
Dispersion & 1.000 & 0 (ex.) & 0 (ex.) & 0 (ex.) & \textbf{+0.021 [+0.01, +0.03]} & +0.000 [$-$0.00, +0.00] & $-$0.000 [$-$0.00, +0.00] & +0.000 [$-$0.00, +0.00] & +0.000 [$-$0.00, +0.00] \\
Classic & 0.969 & \textbf{$-$0.022 [$-$0.03, $-$0.01]} & 0 (ex.) & \textbf{+0.035 [+0.02, +0.05]} & 0 (ex.) & \textbf{$-$0.003 [$-$0.01, $-$0.00]} & \textbf{$-$0.006 [$-$0.01, $-$0.00]} & \textbf{$-$0.007 [$-$0.01, $-$0.00]} & $-$0.004 [$-$0.01, +0.00] \\
Dur.\ trend & 0.957 & +0.003 [$-$0.00, +0.01] & +0.003 [$-$0.00, +0.01] & \textbf{+0.473 [+0.44, +0.50]} & 0 (ex.) & 0 (ex.) & 0 (ex.) & 0 (ex.) & 0 (ex.) \\
RQA $0.5u$ & 0.999 & \textbf{+0.010 [+0.00, +0.02]} & 0 (ex.) & 0 (ex.) & 0 (ex.) & +0.001 [$-$0.00, +0.00] & +0.002 [$-$0.00, +0.00] & \textbf{+0.004 [+0.00, +0.01]} & \textbf{+0.015 [+0.01, +0.02]} \\
RQA $1u$ & 1.000 & \textbf{+0.006 [+0.00, +0.01]} & 0 (ex.) & 0 (ex.) & 0 (ex.) & +0.000 [$-$0.00, +0.00] & +0.001 [$-$0.00, +0.00] & +0.001 [$-$0.00, +0.00] & \textbf{+0.006 [+0.00, +0.01]} \\
Directed & 0.977 & \textbf{$-$0.014 [$-$0.02, $-$0.01]} & 0 (ex.) & 0 (ex.) & \textbf{$-$0.007 [$-$0.01, $-$0.00]} & \textbf{$-$0.011 [$-$0.02, $-$0.00]} & \textbf{$-$0.015 [$-$0.02, $-$0.01]} & \textbf{$-$0.016 [$-$0.02, $-$0.01]} & \textbf{$-$0.013 [$-$0.02, $-$0.01]} \\
FAG lattice & 0.988 & 0 (ex.) & $-$0.001 [$-$0.00, +0.00] & +0.004 [$-$0.00, +0.01] & +0.002 [$-$0.00, +0.01] & $-$0.000 [$-$0.01, +0.00] & +0.001 [$-$0.00, +0.01] & \textbf{+0.009 [+0.00, +0.02]} & \textbf{+0.049 [+0.04, +0.06]} \\
\bottomrule
\end{tabular}}
\end{table}

\begin{table}[!htbp]\centering\scriptsize
\caption{Profile of \texttt{gb\_vd1\_blg}: video viewing vs. ball-catching game. Sensitivity (AUC drop) with refit-bootstrap 95\% interval; bold: interval excludes zero.}
\resizebox{\linewidth}{!}{\begin{tabular}{lrllllllll}
\toprule
Descriptor & AUC & order & dur.\ coupl. & dur.\ marg. & rotation & jit.\ 0.25$u$ & jit.\ 0.5$u$ & jit.\ 1$u$ & jit.\ 2$u$ \\
\midrule
Count & 0.690 & 0 (ex.) & 0 (ex.) & 0 (ex.) & 0 (ex.) & 0 (ex.) & 0 (ex.) & 0 (ex.) & 0 (ex.) \\
Dispersion & 1.000 & 0 (ex.) & 0 (ex.) & 0 (ex.) & +0.000 [$-$0.00, +0.00] & +0.000 [$-$0.00, +0.00] & +0.000 [$-$0.00, +0.00] & $-$0.000 [$-$0.00, +0.00] & +0.000 [$-$0.00, +0.00] \\
Classic & 0.886 & \textbf{$-$0.114 [$-$0.14, $-$0.09]} & 0 (ex.) & \textbf{+0.033 [+0.02, +0.05]} & 0 (ex.) & \textbf{$-$0.003 [$-$0.00, $-$0.00]} & \textbf{$-$0.007 [$-$0.01, $-$0.00]} & \textbf{$-$0.014 [$-$0.02, $-$0.01]} & \textbf{$-$0.014 [$-$0.02, $-$0.01]} \\
Dur.\ trend & 0.752 & +0.011 [$-$0.00, +0.02] & +0.011 [$-$0.00, +0.02] & \textbf{+0.252 [+0.20, +0.30]} & 0 (ex.) & 0 (ex.) & 0 (ex.) & 0 (ex.) & 0 (ex.) \\
RQA $0.5u$ & 0.999 & \textbf{+0.013 [+0.01, +0.02]} & 0 (ex.) & 0 (ex.) & 0 (ex.) & $-$0.000 [$-$0.00, +0.00] & $-$0.000 [$-$0.00, +0.00] & $-$0.000 [$-$0.00, +0.00] & \textbf{+0.003 [+0.00, +0.01]} \\
RQA $1u$ & 1.000 & \textbf{+0.001 [+0.00, +0.00]} & 0 (ex.) & 0 (ex.) & 0 (ex.) & $-$0.000 [$-$0.00, +0.00] & $-$0.000 [$-$0.00, +0.00] & +0.000 [$-$0.00, +0.00] & +0.000 [$-$0.00, +0.00] \\
Directed & 1.000 & \textbf{+0.005 [+0.00, +0.01]} & 0 (ex.) & 0 (ex.) & +0.000 [$-$0.00, +0.00] & +0.000 [$-$0.00, +0.00] & +0.000 [$-$0.00, +0.00] & +0.000 [$-$0.00, +0.00] & \textbf{+0.002 [+0.00, +0.00]} \\
FAG lattice & 0.957 & 0 (ex.) & +0.000 [$-$0.00, +0.00] & \textbf{$-$0.008 [$-$0.01, $-$0.00]} & \textbf{+0.029 [+0.02, +0.04]} & +0.007 [$-$0.00, +0.02] & +0.010 [$-$0.00, +0.02] & \textbf{+0.019 [+0.00, +0.04]} & \textbf{+0.027 [+0.01, +0.04]} \\
\bottomrule
\end{tabular}}
\end{table}

\begin{table}[!htbp]\centering\scriptsize
\caption{Profile of \texttt{gb\_vd1\_vd2}: video segment 1 vs. 2. Sensitivity (AUC drop) with refit-bootstrap 95\% interval; bold: interval excludes zero.}
\resizebox{\linewidth}{!}{\begin{tabular}{lrllllllll}
\toprule
Descriptor & AUC & order & dur.\ coupl. & dur.\ marg. & rotation & jit.\ 0.25$u$ & jit.\ 0.5$u$ & jit.\ 1$u$ & jit.\ 2$u$ \\
\midrule
Count & 0.634 & 0 (ex.) & 0 (ex.) & 0 (ex.) & 0 (ex.) & 0 (ex.) & 0 (ex.) & 0 (ex.) & 0 (ex.) \\
Dispersion & 0.707 & 0 (ex.) & 0 (ex.) & 0 (ex.) & \textbf{+0.098 [+0.07, +0.13]} & +0.004 [$-$0.00, +0.01] & \textbf{+0.012 [+0.00, +0.02]} & \textbf{+0.045 [+0.03, +0.06]} & \textbf{+0.116 [+0.08, +0.15]} \\
Classic & 0.640 & +0.007 [$-$0.00, +0.01] & 0 (ex.) & \textbf{+0.124 [+0.10, +0.15]} & 0 (ex.) & +0.001 [$-$0.00, +0.00] & \textbf{+0.004 [+0.00, +0.01]} & +0.005 [$-$0.00, +0.01] & +0.005 [$-$0.00, +0.01] \\
Dur.\ trend & 0.790 & \textbf{+0.157 [+0.13, +0.19]} & \textbf{+0.158 [+0.13, +0.19]} & \textbf{+0.286 [+0.23, +0.34]} & 0 (ex.) & 0 (ex.) & 0 (ex.) & 0 (ex.) & 0 (ex.) \\
RQA $0.5u$ & 0.677 & \textbf{+0.035 [+0.01, +0.06]} & 0 (ex.) & 0 (ex.) & 0 (ex.) & $-$0.006 [$-$0.03, +0.01] & +0.002 [$-$0.03, +0.03] & \textbf{+0.046 [+0.00, +0.09]} & \textbf{+0.108 [+0.07, +0.14]} \\
RQA $1u$ & 0.719 & \textbf{+0.073 [+0.04, +0.10]} & 0 (ex.) & 0 (ex.) & 0 (ex.) & +0.007 [$-$0.01, +0.02] & +0.013 [$-$0.01, +0.03] & \textbf{+0.043 [+0.01, +0.08]} & \textbf{+0.118 [+0.08, +0.16]} \\
Directed & 0.700 & \textbf{+0.040 [+0.02, +0.06]} & 0 (ex.) & 0 (ex.) & $-$0.010 [$-$0.03, +0.01] & $-$0.005 [$-$0.03, +0.02] & +0.013 [$-$0.01, +0.04] & \textbf{+0.042 [+0.01, +0.08]} & \textbf{+0.069 [+0.04, +0.10]} \\
FAG lattice & 0.664 & 0 (ex.) & \textbf{$-$0.015 [$-$0.03, $-$0.00]} & \textbf{$-$0.020 [$-$0.03, $-$0.01]} & $-$0.012 [$-$0.05, +0.02] & $-$0.015 [$-$0.05, +0.02] & $-$0.006 [$-$0.05, +0.04] & +0.017 [$-$0.03, +0.06] & \textbf{+0.055 [+0.01, +0.10]} \\
\bottomrule
\end{tabular}}
\end{table}

\begin{table}[!htbp]\centering\scriptsize
\caption{Profile of \texttt{gb\_vd1\_halves}: video 1, first vs. second half. Sensitivity (AUC drop) with refit-bootstrap 95\% interval; bold: interval excludes zero.}
\resizebox{\linewidth}{!}{\begin{tabular}{lrllllllll}
\toprule
Descriptor & AUC & order & dur.\ coupl. & dur.\ marg. & rotation & jit.\ 0.25$u$ & jit.\ 0.5$u$ & jit.\ 1$u$ & jit.\ 2$u$ \\
\midrule
Count & 0.500 & 0 (ex.) & 0 (ex.) & 0 (ex.) & 0 (ex.) & 0 (ex.) & 0 (ex.) & 0 (ex.) & 0 (ex.) \\
Dispersion & 0.924 & 0 (ex.) & 0 (ex.) & 0 (ex.) & \textbf{+0.249 [+0.22, +0.27]} & +0.001 [$-$0.00, +0.00] & \textbf{+0.007 [+0.00, +0.01]} & \textbf{+0.029 [+0.02, +0.04]} & \textbf{+0.130 [+0.11, +0.15]} \\
Classic & 0.564 & \textbf{$-$0.090 [$-$0.11, $-$0.07]} & 0 (ex.) & +0.008 [$-$0.01, +0.03] & 0 (ex.) & $-$0.001 [$-$0.00, +0.00] & +0.002 [$-$0.00, +0.01] & +0.005 [$-$0.01, +0.02] & +0.015 [$-$0.00, +0.03] \\
Dur.\ trend & 0.590 & \textbf{+0.076 [+0.03, +0.12]} & \textbf{+0.085 [+0.04, +0.13]} & \textbf{+0.082 [+0.04, +0.12]} & 0 (ex.) & 0 (ex.) & 0 (ex.) & 0 (ex.) & 0 (ex.) \\
RQA $0.5u$ & 0.743 & \textbf{+0.046 [+0.02, +0.07]} & 0 (ex.) & 0 (ex.) & 0 (ex.) & +0.012 [$-$0.01, +0.03] & +0.025 [$-$0.00, +0.05] & \textbf{+0.079 [+0.05, +0.11]} & \textbf{+0.143 [+0.10, +0.19]} \\
RQA $1u$ & 0.823 & \textbf{+0.088 [+0.06, +0.11]} & 0 (ex.) & 0 (ex.) & 0 (ex.) & +0.003 [$-$0.01, +0.02] & \textbf{+0.026 [+0.01, +0.05]} & \textbf{+0.088 [+0.05, +0.12]} & \textbf{+0.182 [+0.15, +0.21]} \\
Directed & 0.805 & \textbf{+0.036 [+0.02, +0.05]} & 0 (ex.) & 0 (ex.) & $-$0.007 [$-$0.03, +0.01] & +0.010 [$-$0.01, +0.03] & \textbf{+0.033 [+0.01, +0.05]} & \textbf{+0.090 [+0.06, +0.12]} & \textbf{+0.191 [+0.16, +0.22]} \\
FAG lattice & 0.741 & 0 (ex.) & $-$0.003 [$-$0.01, +0.00] & +0.000 [$-$0.00, +0.01] & +0.023 [$-$0.01, +0.05] & +0.004 [$-$0.02, +0.03] & +0.011 [$-$0.02, +0.05] & +0.037 [$-$0.00, +0.07] & \textbf{+0.088 [+0.05, +0.12]} \\
\bottomrule
\end{tabular}}
\end{table}

\begin{table}[!htbp]\centering\scriptsize
\caption{Profile of \texttt{gb\_tex\_s1s2}: reading, session 1 vs. 2. Sensitivity (AUC drop) with refit-bootstrap 95\% interval; bold: interval excludes zero.}
\resizebox{\linewidth}{!}{\begin{tabular}{lrllllllll}
\toprule
Descriptor & AUC & order & dur.\ coupl. & dur.\ marg. & rotation & jit.\ 0.25$u$ & jit.\ 0.5$u$ & jit.\ 1$u$ & jit.\ 2$u$ \\
\midrule
Count & 0.530 & 0 (ex.) & 0 (ex.) & 0 (ex.) & 0 (ex.) & 0 (ex.) & 0 (ex.) & 0 (ex.) & 0 (ex.) \\
Dispersion & 0.519 & 0 (ex.) & 0 (ex.) & 0 (ex.) & +0.036 [$-$0.05, +0.12] & +0.003 [$-$0.01, +0.01] & $-$0.000 [$-$0.03, +0.03] & +0.011 [$-$0.03, +0.06] & +0.031 [$-$0.01, +0.07] \\
Classic & 0.539 & +0.005 [$-$0.02, +0.03] & 0 (ex.) & +0.031 [$-$0.04, +0.10] & 0 (ex.) & +0.001 [$-$0.00, +0.01] & $-$0.000 [$-$0.01, +0.01] & $-$0.005 [$-$0.02, +0.02] & $-$0.001 [$-$0.03, +0.03] \\
Dur.\ trend & 0.549 & +0.014 [$-$0.04, +0.07] & +0.016 [$-$0.03, +0.06] & +0.036 [$-$0.05, +0.12] & 0 (ex.) & 0 (ex.) & 0 (ex.) & 0 (ex.) & 0 (ex.) \\
RQA $0.5u$ & 0.576 & \textbf{+0.093 [+0.01, +0.17]} & 0 (ex.) & 0 (ex.) & 0 (ex.) & +0.037 [$-$0.05, +0.12] & \textbf{+0.071 [+0.02, +0.12]} & +0.072 [$-$0.01, +0.15] & +0.062 [$-$0.03, +0.16] \\
RQA $1u$ & 0.521 & +0.047 [$-$0.03, +0.12] & 0 (ex.) & 0 (ex.) & 0 (ex.) & $-$0.007 [$-$0.05, +0.04] & +0.017 [$-$0.06, +0.09] & +0.028 [$-$0.06, +0.12] & +0.007 [$-$0.08, +0.09] \\
Directed & 0.531 & +0.028 [$-$0.05, +0.11] & 0 (ex.) & 0 (ex.) & +0.007 [$-$0.06, +0.07] & +0.003 [$-$0.06, +0.06] & +0.020 [$-$0.03, +0.07] & +0.038 [$-$0.02, +0.10] & +0.039 [$-$0.03, +0.11] \\
FAG lattice & 0.591 & 0 (ex.) & +0.010 [$-$0.01, +0.03] & +0.014 [$-$0.01, +0.04] & +0.042 [$-$0.03, +0.11] & +0.036 [$-$0.04, +0.11] & +0.056 [$-$0.03, +0.14] & +0.067 [$-$0.02, +0.16] & +0.081 [$-$0.01, +0.17] \\
\bottomrule
\end{tabular}}
\end{table}

\begin{table}[!htbp]\centering\scriptsize
\caption{Profile of \texttt{sb\_native}: native vs. non-native English readers, reading pages. Sensitivity (AUC drop) with refit-bootstrap 95\% interval; bold: interval excludes zero.}
\resizebox{\linewidth}{!}{\begin{tabular}{lrllllllll}
\toprule
Descriptor & AUC & order & dur.\ coupl. & dur.\ marg. & rotation & jit.\ 0.25$u$ & jit.\ 0.5$u$ & jit.\ 1$u$ & jit.\ 2$u$ \\
\midrule
Count & 0.624 & 0 (ex.) & 0 (ex.) & 0 (ex.) & 0 (ex.) & 0 (ex.) & 0 (ex.) & 0 (ex.) & 0 (ex.) \\
Dispersion & 0.545 & 0 (ex.) & 0 (ex.) & 0 (ex.) & +0.005 [$-$0.06, +0.07] & +0.000 [$-$0.01, +0.01] & +0.000 [$-$0.02, +0.02] & $-$0.003 [$-$0.04, +0.03] & +0.006 [$-$0.06, +0.07] \\
Classic & 0.689 & +0.037 [$-$0.08, +0.15] & 0 (ex.) & +0.017 [$-$0.07, +0.10] & 0 (ex.) & +0.001 [$-$0.00, +0.00] & +0.004 [$-$0.00, +0.01] & +0.012 [$-$0.01, +0.03] & +0.022 [$-$0.04, +0.08] \\
Dur.\ trend & 0.654 & +0.011 [$-$0.04, +0.06] & +0.012 [$-$0.03, +0.06] & \textbf{+0.144 [+0.00, +0.29]} & 0 (ex.) & 0 (ex.) & 0 (ex.) & 0 (ex.) & 0 (ex.) \\
RQA $0.5u$ & 0.614 & +0.074 [$-$0.01, +0.16] & 0 (ex.) & 0 (ex.) & 0 (ex.) & +0.012 [$-$0.02, +0.04] & +0.029 [$-$0.01, +0.07] & \textbf{+0.079 [+0.01, +0.15]} & \textbf{+0.088 [+0.01, +0.16]} \\
RQA $1u$ & 0.669 & +0.105 [$-$0.00, +0.21] & 0 (ex.) & 0 (ex.) & 0 (ex.) & $-$0.006 [$-$0.03, +0.02] & +0.015 [$-$0.02, +0.05] & \textbf{+0.086 [+0.02, +0.16]} & \textbf{+0.140 [+0.05, +0.23]} \\
Directed & 0.614 & $-$0.030 [$-$0.09, +0.03] & 0 (ex.) & 0 (ex.) & $-$0.004 [$-$0.04, +0.03] & $-$0.011 [$-$0.04, +0.02] & $-$0.015 [$-$0.06, +0.03] & $-$0.014 [$-$0.06, +0.04] & $-$0.004 [$-$0.07, +0.06] \\
FAG lattice & 0.587 & 0 (ex.) & $-$0.000 [$-$0.02, +0.02] & +0.008 [$-$0.02, +0.03] & $-$0.002 [$-$0.04, +0.03] & $-$0.004 [$-$0.03, +0.02] & $-$0.011 [$-$0.06, +0.03] & $-$0.018 [$-$0.07, +0.04] & $-$0.022 [$-$0.09, +0.04] \\
\bottomrule
\end{tabular}}
\end{table}

\begin{table}[!htbp]\centering\scriptsize
\caption{Profile of \texttt{sb\_compr}: high vs. low passage comprehension, reading pages. Sensitivity (AUC drop) with refit-bootstrap 95\% interval; bold: interval excludes zero.}
\resizebox{\linewidth}{!}{\begin{tabular}{lrllllllll}
\toprule
Descriptor & AUC & order & dur.\ coupl. & dur.\ marg. & rotation & jit.\ 0.25$u$ & jit.\ 0.5$u$ & jit.\ 1$u$ & jit.\ 2$u$ \\
\midrule
Count & 0.554 & 0 (ex.) & 0 (ex.) & 0 (ex.) & 0 (ex.) & 0 (ex.) & 0 (ex.) & 0 (ex.) & 0 (ex.) \\
Dispersion & 0.413 & 0 (ex.) & 0 (ex.) & 0 (ex.) & $-$0.047 [$-$0.12, +0.03] & $-$0.004 [$-$0.02, +0.01] & $-$0.007 [$-$0.03, +0.01] & \textbf{$-$0.035 [$-$0.06, $-$0.01]} & $-$0.041 [$-$0.11, +0.03] \\
Classic & 0.587 & +0.010 [$-$0.07, +0.09] & 0 (ex.) & +0.026 [$-$0.09, +0.14] & 0 (ex.) & +0.000 [$-$0.00, +0.00] & +0.001 [$-$0.01, +0.01] & +0.000 [$-$0.01, +0.01] & +0.003 [$-$0.04, +0.05] \\
Dur.\ trend & 0.575 & $-$0.006 [$-$0.04, +0.03] & $-$0.007 [$-$0.03, +0.02] & +0.070 [$-$0.05, +0.19] & 0 (ex.) & 0 (ex.) & 0 (ex.) & 0 (ex.) & 0 (ex.) \\
RQA $0.5u$ & 0.497 & +0.011 [$-$0.12, +0.14] & 0 (ex.) & 0 (ex.) & 0 (ex.) & +0.009 [$-$0.06, +0.07] & $-$0.003 [$-$0.10, +0.09] & +0.001 [$-$0.08, +0.08] & $-$0.005 [$-$0.11, +0.09] \\
RQA $1u$ & 0.529 & +0.050 [$-$0.04, +0.14] & 0 (ex.) & 0 (ex.) & 0 (ex.) & $-$0.017 [$-$0.06, +0.03] & $-$0.014 [$-$0.07, +0.04] & +0.006 [$-$0.05, +0.07] & +0.032 [$-$0.08, +0.15] \\
Directed & 0.533 & +0.004 [$-$0.07, +0.08] & 0 (ex.) & 0 (ex.) & $-$0.010 [$-$0.10, +0.07] & +0.013 [$-$0.05, +0.07] & $-$0.004 [$-$0.08, +0.07] & +0.003 [$-$0.07, +0.08] & +0.010 [$-$0.07, +0.09] \\
FAG lattice & 0.506 & 0 (ex.) & $-$0.007 [$-$0.04, +0.02] & $-$0.002 [$-$0.03, +0.03] & $-$0.048 [$-$0.10, +0.00] & $-$0.032 [$-$0.10, +0.04] & $-$0.035 [$-$0.11, +0.04] & $-$0.038 [$-$0.13, +0.05] & $-$0.011 [$-$0.10, +0.08] \\
\bottomrule
\end{tabular}}
\end{table}

\begin{table}[!htbp]\centering\scriptsize
\caption{Profile of \texttt{em\_program}: Rectangle vs. Vehicle program. Sensitivity (AUC drop) with refit-bootstrap 95\% interval; bold: interval excludes zero.}
\resizebox{\linewidth}{!}{\begin{tabular}{lrllllllll}
\toprule
Descriptor & AUC & order & dur.\ coupl. & dur.\ marg. & rotation & jit.\ 0.25$u$ & jit.\ 0.5$u$ & jit.\ 1$u$ & jit.\ 2$u$ \\
\midrule
Count & 0.561 & 0 (ex.) & 0 (ex.) & 0 (ex.) & 0 (ex.) & 0 (ex.) & 0 (ex.) & 0 (ex.) & 0 (ex.) \\
Dispersion & 0.957 & 0 (ex.) & 0 (ex.) & 0 (ex.) & \textbf{+0.139 [+0.08, +0.20]} & $-$0.001 [$-$0.01, +0.00] & $-$0.001 [$-$0.01, +0.01] & +0.001 [$-$0.01, +0.01] & +0.002 [$-$0.01, +0.02] \\
Classic & 0.539 & \textbf{$-$0.249 [$-$0.34, $-$0.16]} & 0 (ex.) & +0.032 [$-$0.08, +0.15] & 0 (ex.) & $-$0.000 [$-$0.02, +0.02] & +0.003 [$-$0.03, +0.04] & $-$0.004 [$-$0.06, +0.05] & +0.008 [$-$0.10, +0.11] \\
Dur.\ trend & 0.540 & +0.017 [$-$0.13, +0.16] & +0.038 [$-$0.12, +0.20] & +0.055 [$-$0.07, +0.18] & 0 (ex.) & 0 (ex.) & 0 (ex.) & 0 (ex.) & 0 (ex.) \\
RQA $0.5u$ & 0.580 & +0.036 [$-$0.07, +0.14] & 0 (ex.) & 0 (ex.) & 0 (ex.) & +0.004 [$-$0.09, +0.09] & +0.061 [$-$0.08, +0.21] & +0.068 [$-$0.12, +0.25] & +0.000 [$-$0.17, +0.17] \\
RQA $1u$ & 0.547 & +0.003 [$-$0.15, +0.16] & 0 (ex.) & 0 (ex.) & 0 (ex.) & $-$0.010 [$-$0.10, +0.08] & +0.003 [$-$0.13, +0.13] & $-$0.046 [$-$0.22, +0.13] & $-$0.067 [$-$0.26, +0.12] \\
Directed & 0.655 & $-$0.003 [$-$0.08, +0.07] & 0 (ex.) & 0 (ex.) & +0.090 [$-$0.00, +0.18] & +0.112 [$-$0.02, +0.24] & +0.141 [$-$0.01, +0.29] & +0.035 [$-$0.08, +0.15] & +0.044 [$-$0.10, +0.19] \\
FAG lattice & 0.658 & 0 (ex.) & $-$0.002 [$-$0.03, +0.03] & +0.001 [$-$0.03, +0.03] & +0.025 [$-$0.12, +0.17] & +0.022 [$-$0.04, +0.08] & $-$0.049 [$-$0.14, +0.04] & $-$0.032 [$-$0.17, +0.11] & +0.005 [$-$0.14, +0.15] \\
FAG AOI & 0.778 & 0 (ex.) & $-$0.005 [$-$0.02, +0.01] & $-$0.005 [$-$0.02, +0.01] & +0.079 [$-$0.04, +0.20] & $-$0.022 [$-$0.08, +0.04] & $-$0.013 [$-$0.10, +0.07] & $-$0.021 [$-$0.14, +0.09] & $-$0.044 [$-$0.14, +0.05] \\
FAG AOI pad. & 0.729 & 0 (ex.) & +0.001 [$-$0.02, +0.03] & +0.001 [$-$0.03, +0.03] & +0.056 [$-$0.06, +0.17] & $-$0.023 [$-$0.10, +0.06] & +0.017 [$-$0.09, +0.12] & +0.028 [$-$0.08, +0.13] & +0.049 [$-$0.09, +0.19] \\
\bottomrule
\end{tabular}}
\end{table}

\begin{table}[!htbp]\centering\scriptsize
\caption{Profile of \texttt{em\_halves}: first vs. second half of a program reading. Sensitivity (AUC drop) with refit-bootstrap 95\% interval; bold: interval excludes zero.}
\resizebox{\linewidth}{!}{\begin{tabular}{lrllllllll}
\toprule
Descriptor & AUC & order & dur.\ coupl. & dur.\ marg. & rotation & jit.\ 0.25$u$ & jit.\ 0.5$u$ & jit.\ 1$u$ & jit.\ 2$u$ \\
\midrule
Count & 0.500 & 0 (ex.) & 0 (ex.) & 0 (ex.) & 0 (ex.) & 0 (ex.) & 0 (ex.) & 0 (ex.) & 0 (ex.) \\
Dispersion & 0.616 & 0 (ex.) & 0 (ex.) & 0 (ex.) & +0.012 [$-$0.05, +0.07] & $-$0.001 [$-$0.01, +0.01] & +0.000 [$-$0.01, +0.01] & +0.001 [$-$0.02, +0.02] & +0.008 [$-$0.02, +0.04] \\
Classic & 0.776 & \textbf{+0.172 [+0.11, +0.23]} & 0 (ex.) & +0.004 [$-$0.01, +0.02] & 0 (ex.) & $-$0.004 [$-$0.01, +0.00] & $-$0.003 [$-$0.02, +0.01] & $-$0.002 [$-$0.03, +0.02] & +0.043 [$-$0.00, +0.09] \\
Dur.\ trend & 0.582 & +0.048 [$-$0.04, +0.14] & +0.036 [$-$0.08, +0.15] & \textbf{+0.086 [+0.00, +0.17]} & 0 (ex.) & 0 (ex.) & 0 (ex.) & 0 (ex.) & 0 (ex.) \\
RQA $0.5u$ & 0.613 & +0.058 [$-$0.04, +0.15] & 0 (ex.) & 0 (ex.) & 0 (ex.) & $-$0.036 [$-$0.09, +0.02] & $-$0.050 [$-$0.11, +0.01] & \textbf{$-$0.085 [$-$0.15, $-$0.02]} & $-$0.053 [$-$0.11, +0.01] \\
RQA $1u$ & 0.680 & +0.081 [$-$0.00, +0.16] & 0 (ex.) & 0 (ex.) & 0 (ex.) & $-$0.006 [$-$0.02, +0.01] & $-$0.018 [$-$0.05, +0.01] & \textbf{$-$0.052 [$-$0.10, $-$0.01]} & $-$0.043 [$-$0.09, +0.01] \\
Directed & 0.634 & +0.033 [$-$0.04, +0.10] & 0 (ex.) & 0 (ex.) & $-$0.004 [$-$0.06, +0.05] & $-$0.010 [$-$0.07, +0.05] & $-$0.010 [$-$0.07, +0.05] & +0.002 [$-$0.07, +0.08] & +0.059 [$-$0.03, +0.15] \\
FAG lattice & 0.599 & 0 (ex.) & $-$0.008 [$-$0.03, +0.01] & $-$0.006 [$-$0.03, +0.01] & $-$0.010 [$-$0.07, +0.05] & $-$0.007 [$-$0.07, +0.05] & $-$0.000 [$-$0.08, +0.07] & +0.006 [$-$0.07, +0.08] & +0.036 [$-$0.08, +0.15] \\
FAG AOI & 0.682 & 0 (ex.) & $-$0.000 [$-$0.01, +0.01] & $-$0.002 [$-$0.01, +0.01] & \textbf{+0.098 [+0.03, +0.17]} & +0.036 [$-$0.03, +0.10] & +0.023 [$-$0.02, +0.07] & +0.045 [$-$0.03, +0.12] & \textbf{+0.113 [+0.01, +0.21]} \\
FAG AOI pad. & 0.669 & 0 (ex.) & $-$0.004 [$-$0.02, +0.02] & $-$0.001 [$-$0.02, +0.02] & +0.047 [$-$0.02, +0.11] & +0.027 [$-$0.04, +0.09] & +0.016 [$-$0.06, +0.09] & +0.016 [$-$0.05, +0.09] & +0.062 [$-$0.02, +0.14] \\
\bottomrule
\end{tabular}}
\end{table}

\begin{table}[!htbp]\centering\scriptsize
\caption{Profile of \texttt{em\_expertise}: high vs. low/no self-rated programming expertise. Sensitivity (AUC drop) with refit-bootstrap 95\% interval; bold: interval excludes zero.}
\resizebox{\linewidth}{!}{\begin{tabular}{lrllllllll}
\toprule
Descriptor & AUC & order & dur.\ coupl. & dur.\ marg. & rotation & jit.\ 0.25$u$ & jit.\ 0.5$u$ & jit.\ 1$u$ & jit.\ 2$u$ \\
\midrule
Count & 0.472 & 0 (ex.) & 0 (ex.) & 0 (ex.) & 0 (ex.) & 0 (ex.) & 0 (ex.) & 0 (ex.) & 0 (ex.) \\
Dispersion & 0.469 & 0 (ex.) & 0 (ex.) & 0 (ex.) & $-$0.028 [$-$0.21, +0.15] & +0.006 [$-$0.03, +0.04] & +0.004 [$-$0.05, +0.06] & +0.007 [$-$0.06, +0.08] & $-$0.004 [$-$0.10, +0.09] \\
Classic & 0.438 & $-$0.032 [$-$0.28, +0.21] & 0 (ex.) & $-$0.023 [$-$0.34, +0.29] & 0 (ex.) & $-$0.002 [$-$0.04, +0.03] & $-$0.008 [$-$0.05, +0.03] & $-$0.043 [$-$0.10, +0.01] & $-$0.031 [$-$0.16, +0.09] \\
Dur.\ trend & 0.541 & +0.046 [$-$0.13, +0.23] & +0.070 [$-$0.23, +0.37] & +0.031 [$-$0.27, +0.33] & 0 (ex.) & 0 (ex.) & 0 (ex.) & 0 (ex.) & 0 (ex.) \\
RQA $0.5u$ & 0.419 & $-$0.016 [$-$0.37, +0.34] & 0 (ex.) & 0 (ex.) & 0 (ex.) & $-$0.022 [$-$0.21, +0.17] & $-$0.033 [$-$0.11, +0.04] & +0.068 [$-$0.13, +0.27] & +0.056 [$-$0.21, +0.32] \\
RQA $1u$ & 0.416 & $-$0.037 [$-$0.37, +0.30] & 0 (ex.) & 0 (ex.) & 0 (ex.) & $-$0.017 [$-$0.13, +0.09] & +0.016 [$-$0.24, +0.28] & +0.028 [$-$0.21, +0.27] & +0.090 [$-$0.24, +0.42] \\
Directed & 0.366 & +0.040 [$-$0.12, +0.20] & 0 (ex.) & 0 (ex.) & $-$0.018 [$-$0.10, +0.06] & +0.044 [$-$0.17, +0.26] & +0.013 [$-$0.25, +0.28] & $-$0.086 [$-$0.31, +0.14] & $-$0.083 [$-$0.32, +0.15] \\
FAG lattice & 0.447 & 0 (ex.) & +0.022 [$-$0.03, +0.08] & +0.028 [$-$0.07, +0.12] & $-$0.026 [$-$0.25, +0.20] & $-$0.035 [$-$0.28, +0.21] & $-$0.065 [$-$0.28, +0.15] & $-$0.062 [$-$0.39, +0.27] & $-$0.058 [$-$0.29, +0.18] \\
FAG AOI & 0.259 & 0 (ex.) & $-$0.009 [$-$0.07, +0.06] & $-$0.017 [$-$0.11, +0.07] & $-$0.231 [$-$0.48, +0.02] & $-$0.128 [$-$0.35, +0.09] & $-$0.143 [$-$0.37, +0.09] & $-$0.196 [$-$0.49, +0.10] & $-$0.110 [$-$0.39, +0.17] \\
FAG AOI pad. & 0.459 & 0 (ex.) & +0.015 [$-$0.05, +0.08] & +0.014 [$-$0.05, +0.07] & $-$0.034 [$-$0.26, +0.20] & +0.062 [$-$0.20, +0.33] & +0.010 [$-$0.25, +0.27] & $-$0.010 [$-$0.27, +0.25] & $-$0.023 [$-$0.30, +0.26] \\
\bottomrule
\end{tabular}}
\end{table}

\begin{table}[!htbp]\centering\scriptsize
\caption{Profile of \texttt{em\_correct}: correct vs. incorrect answer, both programs pooled (confounded). Sensitivity (AUC drop) with refit-bootstrap 95\% interval; bold: interval excludes zero.}
\resizebox{\linewidth}{!}{\begin{tabular}{lrllllllll}
\toprule
Descriptor & AUC & order & dur.\ coupl. & dur.\ marg. & rotation & jit.\ 0.25$u$ & jit.\ 0.5$u$ & jit.\ 1$u$ & jit.\ 2$u$ \\
\midrule
Count & 0.495 & 0 (ex.) & 0 (ex.) & 0 (ex.) & 0 (ex.) & 0 (ex.) & 0 (ex.) & 0 (ex.) & 0 (ex.) \\
Dispersion & 0.716 & 0 (ex.) & 0 (ex.) & 0 (ex.) & \textbf{+0.096 [+0.01, +0.18]} & +0.000 [$-$0.01, +0.01] & +0.003 [$-$0.01, +0.01] & $-$0.001 [$-$0.01, +0.01] & +0.005 [$-$0.03, +0.04] \\
Classic & 0.501 & $-$0.123 [$-$0.27, +0.03] & 0 (ex.) & +0.054 [$-$0.07, +0.18] & 0 (ex.) & $-$0.000 [$-$0.02, +0.02] & $-$0.000 [$-$0.03, +0.03] & $-$0.004 [$-$0.06, +0.05] & $-$0.002 [$-$0.07, +0.07] \\
Dur.\ trend & 0.511 & $-$0.013 [$-$0.17, +0.15] & $-$0.023 [$-$0.17, +0.13] & +0.046 [$-$0.09, +0.18] & 0 (ex.) & 0 (ex.) & 0 (ex.) & 0 (ex.) & 0 (ex.) \\
RQA $0.5u$ & 0.571 & +0.022 [$-$0.10, +0.14] & 0 (ex.) & 0 (ex.) & 0 (ex.) & +0.040 [$-$0.05, +0.13] & +0.056 [$-$0.05, +0.16] & +0.097 [$-$0.06, +0.25] & +0.061 [$-$0.12, +0.24] \\
RQA $1u$ & 0.525 & $-$0.031 [$-$0.19, +0.13] & 0 (ex.) & 0 (ex.) & 0 (ex.) & +0.018 [$-$0.05, +0.09] & +0.069 [$-$0.04, +0.18] & +0.062 [$-$0.08, +0.20] & +0.009 [$-$0.18, +0.20] \\
Directed & 0.580 & $-$0.021 [$-$0.11, +0.07] & 0 (ex.) & 0 (ex.) & +0.043 [$-$0.05, +0.14] & +0.053 [$-$0.04, +0.15] & +0.103 [$-$0.02, +0.23] & +0.095 [$-$0.05, +0.24] & +0.058 [$-$0.12, +0.23] \\
FAG lattice & 0.664 & 0 (ex.) & +0.001 [$-$0.01, +0.02] & $-$0.002 [$-$0.02, +0.01] & +0.002 [$-$0.07, +0.08] & +0.002 [$-$0.07, +0.07] & $-$0.004 [$-$0.08, +0.08] & +0.051 [$-$0.07, +0.17] & +0.054 [$-$0.07, +0.18] \\
FAG AOI & 0.635 & 0 (ex.) & $-$0.004 [$-$0.03, +0.02] & $-$0.001 [$-$0.02, +0.02] & +0.076 [$-$0.09, +0.24] & +0.014 [$-$0.09, +0.12] & +0.052 [$-$0.06, +0.17] & $-$0.008 [$-$0.15, +0.13] & +0.031 [$-$0.11, +0.17] \\
FAG AOI pad. & 0.571 & 0 (ex.) & $-$0.006 [$-$0.03, +0.02] & $-$0.010 [$-$0.04, +0.02] & +0.052 [$-$0.09, +0.19] & $-$0.020 [$-$0.16, +0.12] & $-$0.022 [$-$0.16, +0.12] & $-$0.007 [$-$0.15, +0.14] & +0.004 [$-$0.15, +0.16] \\
\bottomrule
\end{tabular}}
\end{table}

\begin{table}[!htbp]\centering\scriptsize
\caption{Profile of \texttt{em\_correctrect}: correct vs. incorrect answer, Rectangle program only. Sensitivity (AUC drop) with refit-bootstrap 95\% interval; bold: interval excludes zero.}
\resizebox{\linewidth}{!}{\begin{tabular}{lrllllllll}
\toprule
Descriptor & AUC & order & dur.\ coupl. & dur.\ marg. & rotation & jit.\ 0.25$u$ & jit.\ 0.5$u$ & jit.\ 1$u$ & jit.\ 2$u$ \\
\midrule
Count & 0.589 & 0 (ex.) & 0 (ex.) & 0 (ex.) & 0 (ex.) & 0 (ex.) & 0 (ex.) & 0 (ex.) & 0 (ex.) \\
Dispersion & 0.472 & 0 (ex.) & 0 (ex.) & 0 (ex.) & $-$0.062 [$-$0.25, +0.13] & +0.002 [$-$0.02, +0.02] & +0.006 [$-$0.03, +0.04] & +0.034 [$-$0.02, +0.09] & +0.073 [$-$0.02, +0.17] \\
Classic & 0.523 & $-$0.002 [$-$0.19, +0.19] & 0 (ex.) & +0.048 [$-$0.19, +0.28] & 0 (ex.) & +0.004 [$-$0.03, +0.04] & +0.008 [$-$0.05, +0.06] & $-$0.006 [$-$0.08, +0.06] & +0.003 [$-$0.08, +0.09] \\
Dur.\ trend & 0.507 & $-$0.069 [$-$0.23, +0.09] & $-$0.077 [$-$0.24, +0.08] & $-$0.037 [$-$0.29, +0.22] & 0 (ex.) & 0 (ex.) & 0 (ex.) & 0 (ex.) & 0 (ex.) \\
RQA $0.5u$ & 0.451 & $-$0.054 [$-$0.38, +0.27] & 0 (ex.) & 0 (ex.) & 0 (ex.) & $-$0.046 [$-$0.25, +0.16] & $-$0.072 [$-$0.22, +0.08] & $-$0.004 [$-$0.26, +0.25] & $-$0.041 [$-$0.28, +0.20] \\
RQA $1u$ & 0.462 & $-$0.065 [$-$0.38, +0.25] & 0 (ex.) & 0 (ex.) & 0 (ex.) & $-$0.015 [$-$0.15, +0.12] & $-$0.026 [$-$0.18, +0.13] & $-$0.046 [$-$0.24, +0.15] & $-$0.048 [$-$0.29, +0.20] \\
Directed & 0.533 & +0.006 [$-$0.22, +0.23] & 0 (ex.) & 0 (ex.) & $-$0.013 [$-$0.19, +0.17] & +0.000 [$-$0.18, +0.19] & +0.014 [$-$0.19, +0.22] & $-$0.000 [$-$0.22, +0.22] & +0.018 [$-$0.19, +0.23] \\
FAG lattice & 0.578 & 0 (ex.) & +0.007 [$-$0.05, +0.06] & +0.011 [$-$0.04, +0.06] & +0.021 [$-$0.11, +0.16] & +0.040 [$-$0.15, +0.23] & +0.025 [$-$0.11, +0.16] & +0.057 [$-$0.18, +0.29] & +0.097 [$-$0.06, +0.25] \\
FAG AOI & 0.658 & 0 (ex.) & +0.008 [$-$0.05, +0.07] & +0.005 [$-$0.06, +0.07] & +0.186 [$-$0.03, +0.40] & +0.092 [$-$0.16, +0.34] & \textbf{+0.167 [+0.01, +0.32]} & +0.144 [$-$0.08, +0.37] & +0.201 [$-$0.02, +0.42] \\
FAG AOI pad. & 0.544 & 0 (ex.) & +0.024 [$-$0.03, +0.08] & +0.020 [$-$0.05, +0.09] & +0.029 [$-$0.17, +0.23] & $-$0.027 [$-$0.19, +0.13] & $-$0.003 [$-$0.23, +0.22] & +0.069 [$-$0.18, +0.32] & +0.044 [$-$0.19, +0.28] \\
\bottomrule
\end{tabular}}
\end{table}

\begin{table}[!htbp]\centering\scriptsize
\caption{Profile of \texttt{em\_correctveh}: correct vs. incorrect answer, Vehicle program only. Sensitivity (AUC drop) with refit-bootstrap 95\% interval; bold: interval excludes zero.}
\resizebox{\linewidth}{!}{\begin{tabular}{lrllllllll}
\toprule
Descriptor & AUC & order & dur.\ coupl. & dur.\ marg. & rotation & jit.\ 0.25$u$ & jit.\ 0.5$u$ & jit.\ 1$u$ & jit.\ 2$u$ \\
\midrule
Count & 0.444 & 0 (ex.) & 0 (ex.) & 0 (ex.) & 0 (ex.) & 0 (ex.) & 0 (ex.) & 0 (ex.) & 0 (ex.) \\
Dispersion & 0.476 & 0 (ex.) & 0 (ex.) & 0 (ex.) & +0.091 [$-$0.20, +0.38] & +0.006 [$-$0.02, +0.03] & +0.007 [$-$0.04, +0.06] & +0.015 [$-$0.04, +0.07] & +0.022 [$-$0.10, +0.15] \\
Classic & 0.500 & +0.070 [$-$0.15, +0.29] & 0 (ex.) & +0.025 [$-$0.20, +0.25] & 0 (ex.) & +0.000 [$-$0.04, +0.04] & +0.001 [$-$0.06, +0.06] & +0.003 [$-$0.05, +0.06] & +0.017 [$-$0.10, +0.14] \\
Dur.\ trend & 0.521 & +0.012 [$-$0.22, +0.24] & $-$0.024 [$-$0.26, +0.21] & +0.002 [$-$0.30, +0.31] & 0 (ex.) & 0 (ex.) & 0 (ex.) & 0 (ex.) & 0 (ex.) \\
RQA $0.5u$ & 0.518 & $-$0.026 [$-$0.22, +0.17] & 0 (ex.) & 0 (ex.) & 0 (ex.) & +0.034 [$-$0.15, +0.21] & +0.026 [$-$0.23, +0.28] & +0.104 [$-$0.14, +0.35] & +0.098 [$-$0.11, +0.31] \\
RQA $1u$ & 0.559 & +0.036 [$-$0.24, +0.31] & 0 (ex.) & 0 (ex.) & 0 (ex.) & $-$0.004 [$-$0.11, +0.10] & +0.093 [$-$0.10, +0.29] & +0.106 [$-$0.09, +0.30] & +0.140 [$-$0.07, +0.35] \\
Directed & 0.618 & +0.134 [$-$0.06, +0.33] & 0 (ex.) & 0 (ex.) & +0.080 [$-$0.14, +0.31] & +0.097 [$-$0.11, +0.30] & +0.120 [$-$0.05, +0.29] & +0.147 [$-$0.07, +0.37] & +0.088 [$-$0.15, +0.33] \\
FAG lattice & 0.574 & 0 (ex.) & +0.023 [$-$0.02, +0.07] & +0.021 [$-$0.05, +0.09] & +0.074 [$-$0.18, +0.33] & +0.066 [$-$0.13, +0.26] & +0.032 [$-$0.20, +0.26] & +0.069 [$-$0.17, +0.31] & +0.054 [$-$0.18, +0.29] \\
FAG AOI & 0.503 & 0 (ex.) & +0.000 [$-$0.06, +0.06] & $-$0.004 [$-$0.07, +0.06] & $-$0.010 [$-$0.29, +0.27] & $-$0.011 [$-$0.24, +0.22] & $-$0.017 [$-$0.30, +0.26] & +0.014 [$-$0.20, +0.23] & +0.001 [$-$0.26, +0.27] \\
FAG AOI pad. & 0.470 & 0 (ex.) & $-$0.007 [$-$0.05, +0.03] & $-$0.007 [$-$0.05, +0.04] & $-$0.148 [$-$0.43, +0.13] & +0.010 [$-$0.22, +0.24] & $-$0.052 [$-$0.26, +0.16] & +0.003 [$-$0.24, +0.24] & +0.015 [$-$0.24, +0.27] \\
\bottomrule
\end{tabular}}
\end{table}


\begin{table}[!htbp]\centering\scriptsize
\caption{Observed pooled AUC (all out-of-fold decision values ranked together; class-balanced participant folds), with the permutation $p$-value in brackets (200 permutations under the schemes of Table 2 of the article).}
\resizebox{\linewidth}{!}{\begin{tabular}{lrrrrrrrrrr}
\toprule
Problem & Count & Dispersion & Classic & Dur.\ trend & RQA $0.5u$ & RQA $1u$ & Directed & FAG lattice & FAG AOI & FAG AOI pad. \\
\midrule
\texttt{gb\_tex\_vd1} & 0.968 (.005) & 1.000 (.005) & 0.968 (.005) & 0.956 (.005) & 0.999 (.005) & 1.000 (.005) & 0.977 (.005) & 0.989 (.005) & -- & -- \\
\texttt{gb\_vd1\_blg} & 0.691 (.005) & 1.000 (.005) & 0.888 (.005) & 0.751 (.005) & 0.999 (.005) & 1.000 (.005) & 1.000 (.005) & 0.957 (.005) & -- & -- \\
\texttt{gb\_vd1\_vd2} & 0.633 (.005) & 0.703 (.005) & 0.636 (.005) & 0.791 (.005) & 0.672 (.005) & 0.715 (.005) & 0.700 (.005) & 0.666 (.005) & -- & -- \\
\texttt{gb\_vd1\_halves} & 0.500 (1.000) & 0.921 (.005) & 0.561 (.005) & 0.590 (.005) & 0.742 (.005) & 0.822 (.005) & 0.803 (.005) & 0.739 (.005) & -- & -- \\
\texttt{gb\_tex\_s1s2} & 0.516 (.065) & 0.517 (.224) & 0.537 (.005) & 0.545 (.015) & 0.574 (.005) & 0.515 (.159) & 0.531 (.070) & 0.588 (.005) & -- & -- \\
\texttt{sb\_native} & 0.610 (.010) & 0.530 (.164) & 0.673 (.005) & 0.637 (.015) & 0.611 (.005) & 0.659 (.005) & 0.594 (.025) & 0.575 (.050) & -- & -- \\
\texttt{sb\_compr} & 0.527 (.915) & 0.402 (.990) & 0.595 (.950) & 0.575 (.910) & 0.485 (.896) & 0.523 (.891) & 0.514 (.935) & 0.473 (.955) & -- & -- \\
\texttt{em\_program} & 0.547 (.025) & 0.963 (.005) & 0.546 (.025) & 0.524 (.249) & 0.587 (.015) & 0.541 (.129) & 0.653 (.005) & 0.666 (.005) & 0.781 (.005) & 0.724 (.005) \\
\texttt{em\_halves} & 0.500 (1.000) & 0.619 (.005) & 0.768 (.005) & 0.574 (.010) & 0.606 (.005) & 0.677 (.005) & 0.629 (.005) & 0.602 (.005) & 0.675 (.005) & 0.664 (.005) \\
\texttt{em\_expertise} & 0.400 (.697) & 0.449 (.637) & 0.378 (.816) & 0.510 (.373) & 0.425 (.771) & 0.371 (.881) & 0.319 (.940) & 0.414 (.816) & 0.264 (.985) & 0.418 (.756) \\
\texttt{em\_correct} & 0.463 (.741) & 0.710 (.478) & 0.497 (.522) & 0.505 (.393) & 0.554 (.179) & 0.518 (.328) & 0.559 (.279) & 0.659 (.015) & 0.638 (.254) & 0.569 (.746) \\
\texttt{em\_correctrect} & 0.539 (.214) & 0.446 (.741) & 0.488 (.557) & 0.530 (.303) & 0.455 (.692) & 0.452 (.726) & 0.518 (.368) & 0.556 (.184) & 0.634 (.035) & 0.514 (.403) \\
\texttt{em\_correctveh} & 0.434 (.701) & 0.445 (.731) & 0.458 (.672) & 0.474 (.627) & 0.490 (.507) & 0.564 (.159) & 0.588 (.119) & 0.536 (.299) & 0.465 (.592) & 0.465 (.627) \\
\bottomrule
\end{tabular}
}
\end{table}

\begin{table}[!htbp]\centering\scriptsize
\caption{Range, over descriptors, of the mean of the permutation distribution of the observed AUC, under the within-fold AUC (1000 permutations) and the pooled AUC (200 permutations), both with class-balanced folds. On \texttt{em\_correct} labels are permuted within programs and on \texttt{sb\_compr} between the passages of a participant, so stimulus structure that survives the permutation raises the mean for descriptors that read it.}
\begin{tabular}{lrr}
\toprule
Problem & Within-fold AUC & Pooled AUC \\
\midrule
\texttt{gb\_tex\_vd1} & 0.499--0.501 & 0.491--0.496 \\
\texttt{gb\_vd1\_blg} & 0.499--0.502 & 0.494--0.497 \\
\texttt{gb\_vd1\_vd2} & 0.499--0.501 & 0.494--0.499 \\
\texttt{gb\_vd1\_halves} & 0.499--0.501 & 0.496--0.500 \\
\texttt{gb\_tex\_s1s2} & 0.499--0.501 & 0.494--0.498 \\
\texttt{sb\_native} & 0.497--0.504 & 0.476--0.490 \\
\texttt{sb\_compr} & 0.477--0.618 & 0.458--0.611 \\
\texttt{em\_program} & 0.496--0.501 & 0.488--0.500 \\
\texttt{em\_halves} & 0.498--0.502 & 0.487--0.500 \\
\texttt{em\_expertise} & 0.489--0.504 & 0.466--0.492 \\
\texttt{em\_correct} & 0.502--0.710 & 0.487--0.708 \\
\texttt{em\_correctrect} & 0.496--0.505 & 0.476--0.492 \\
\texttt{em\_correctveh} & 0.496--0.503 & 0.475--0.491 \\
\bottomrule
\end{tabular}

\end{table}

{\scriptsize\setlength{\tabcolsep}{2.5pt}\setlength{\LTcapwidth}{\linewidth}
\begin{longtable}{lllllllllrc}
\caption{Evidence ledger of the 89 interpreted cells. Primary: sensitivity with the within-fold AUC and class-balanced folds ($R = 6$, $B = 60$, seed 42), with its 95\% refit-bootstrap interval (excluding zero by selection). Pooled AUC: the same cell with the pooled AUC. Partitions: range of the point estimate over 20 fold partitions. Long runs: range of the point estimate over three runs with $R = 20$ and $B = 300$ (seeds 42, 43, 44), and the number of those runs whose interval excluded zero. Trees: gradient-boosted trees ($R = 6$, $B = 60$, seed 42). Norm.: sensitivity divided by the observed AUC minus the mean of its permutation distribution, under the within-fold and the pooled AUC. Robust: sign and exclusion of zero hold in the primary run, the pooled run, every partition and every long run and, for the AOI graph, on the padded graph (Table S23 and Section 5.9 of the article). Bold: interval excludes zero.}\\
\toprule
Problem & Descriptor & Transformation & Primary & Pooled AUC & Partitions & Long runs & Trees & Norm. & Norm.\ (pooled) & Robust \\
\midrule
\endfirsthead
\toprule
Problem & Descriptor & Transformation & Primary & Pooled AUC & Partitions & Long runs & Trees & Norm. & Norm.\ (pooled) & Robust \\
\midrule
\endhead
\bottomrule
\endfoot
\texttt{gb\_tex\_vd1} & Classic & duration marginal & +0.035 [+0.021, +0.048] & \textbf{+0.035 [+0.022, +0.048]} & $+0.033$ to $+0.036$ & $+0.034$ to $+0.034$ (3/3) & \textbf{+0.045 [+0.025, +0.065]} & $+0.07$ & $+0.07$ & $\checkmark$ \\
\texttt{gb\_tex\_vd1} & Classic & order & $-$0.022 [$-$0.034, $-$0.010] & \textbf{$-$0.022 [$-$0.034, $-$0.009]} & $-0.024$ to $-0.020$ & $-0.022$ to $-0.022$ (3/3) & \textbf{$-$0.019 [$-$0.033, $-$0.005]} & $-0.05$ & $-0.05$ & $\checkmark$ \\
\texttt{gb\_tex\_vd1} & Classic & jitter $0.25u$ & $-$0.003 [$-$0.005, $-$0.001] & \textbf{$-$0.003 [$-$0.005, $-$0.001]} & $-0.003$ to $-0.002$ & $-0.003$ to $-0.003$ (3/3) & $-$0.002 [$-$0.009, +0.004] & $-0.01$ & $-0.01$ & $\checkmark$ \\
\texttt{gb\_tex\_vd1} & Classic & jitter $0.5u$ & $-$0.006 [$-$0.009, $-$0.002] & \textbf{$-$0.006 [$-$0.008, $-$0.003]} & $-0.007$ to $-0.005$ & $-0.006$ to $-0.005$ (3/3) & $-$0.005 [$-$0.011, +0.001] & $-0.01$ & $-0.01$ & $\checkmark$ \\
\texttt{gb\_tex\_vd1} & Classic & jitter $1u$ & $-$0.007 [$-$0.011, $-$0.003] & \textbf{$-$0.007 [$-$0.011, $-$0.003]} & $-0.009$ to $-0.007$ & $-0.007$ to $-0.007$ (3/3) & $-$0.006 [$-$0.014, +0.003] & $-0.01$ & $-0.02$ & $\checkmark$ \\
\texttt{gb\_tex\_vd1} & Directed & order & $-$0.014 [$-$0.020, $-$0.008] & \textbf{$-$0.014 [$-$0.020, $-$0.008]} & $-0.014$ to $-0.012$ & $-0.015$ to $-0.014$ (3/3) & $-$0.004 [$-$0.011, +0.002] & $-0.03$ & $-0.03$ & $\checkmark$ \\
\texttt{gb\_tex\_vd1} & Directed & jitter $0.25u$ & $-$0.011 [$-$0.018, $-$0.004] & \textbf{$-$0.011 [$-$0.016, $-$0.005]} & $-0.013$ to $-0.011$ & $-0.012$ to $-0.011$ (3/3) & $-$0.004 [$-$0.010, +0.002] & $-0.02$ & $-0.02$ & $\checkmark$ \\
\texttt{gb\_tex\_vd1} & Directed & jitter $0.5u$ & $-$0.015 [$-$0.021, $-$0.009] & \textbf{$-$0.015 [$-$0.020, $-$0.009]} & $-0.015$ to $-0.013$ & $-0.016$ to $-0.015$ (3/3) & $-$0.005 [$-$0.011, +0.001] & $-0.03$ & $-0.03$ & $\checkmark$ \\
\texttt{gb\_tex\_vd1} & Directed & jitter $1u$ & $-$0.016 [$-$0.022, $-$0.011] & \textbf{$-$0.016 [$-$0.022, $-$0.010]} & $-0.017$ to $-0.014$ & $-0.017$ to $-0.016$ (3/3) & $-$0.006 [$-$0.012, +0.000] & $-0.03$ & $-0.03$ & $\checkmark$ \\
\texttt{gb\_tex\_vd1} & Directed & jitter $2u$ & $-$0.013 [$-$0.021, $-$0.005] & \textbf{$-$0.012 [$-$0.019, $-$0.005]} & $-0.013$ to $-0.010$ & $-0.013$ to $-0.012$ (3/3) & $-$0.002 [$-$0.011, +0.007] & $-0.03$ & $-0.03$ & $\checkmark$ \\
\texttt{gb\_tex\_vd1} & Directed & rotation & $-$0.007 [$-$0.013, $-$0.002] & \textbf{$-$0.008 [$-$0.012, $-$0.003]} & $-0.009$ to $-0.007$ & $-0.007$ to $-0.007$ (3/3) & $-$0.002 [$-$0.009, +0.005] & $-0.02$ & $-0.02$ & $\checkmark$ \\
\texttt{gb\_tex\_vd1} & Dispersion & rotation & +0.021 [+0.014, +0.029] & \textbf{+0.021 [+0.014, +0.029]} & $+0.018$ to $+0.020$ & $+0.021$ to $+0.021$ (3/3) & \textbf{+0.016 [+0.008, +0.024]} & $+0.04$ & $+0.04$ & $\checkmark$ \\
\texttt{gb\_tex\_vd1} & Dur.\ trend & duration marginal & +0.473 [+0.442, +0.503] & \textbf{+0.474 [+0.446, +0.503]} & $+0.466$ to $+0.479$ & $+0.454$ to $+0.469$ (3/3) & \textbf{+0.371 [+0.318, +0.424]} & $+1.03$ & $+1.03$ & $\checkmark$ \\
\texttt{gb\_tex\_vd1} & FAG lattice & jitter $1u$ & +0.009 [+0.002, +0.016] & \textbf{+0.010 [+0.002, +0.017]} & $+0.007$ to $+0.009$ & $+0.007$ to $+0.009$ (3/3) & +0.005 [$-$0.003, +0.013] & $+0.02$ & $+0.02$ & $\checkmark$ \\
\texttt{gb\_tex\_vd1} & FAG lattice & jitter $2u$ & +0.049 [+0.037, +0.061] & \textbf{+0.050 [+0.038, +0.062]} & $+0.044$ to $+0.047$ & $+0.047$ to $+0.048$ (3/3) & \textbf{+0.041 [+0.024, +0.058]} & $+0.10$ & $+0.10$ & $\checkmark$ \\
\texttt{gb\_tex\_vd1} & RQA $0.5u$ & order & +0.010 [+0.003, +0.018] & \textbf{+0.011 [+0.004, +0.018]} & $+0.009$ to $+0.011$ & $+0.010$ to $+0.010$ (3/3) & \textbf{+0.013 [+0.004, +0.022]} & $+0.02$ & $+0.02$ & $\checkmark$ \\
\texttt{gb\_tex\_vd1} & RQA $0.5u$ & jitter $1u$ & +0.004 [+0.001, +0.008] & \textbf{+0.004 [+0.001, +0.008]} & $+0.003$ to $+0.004$ & $+0.003$ to $+0.004$ (3/3) & +0.004 [$-$0.000, +0.007] & $+0.01$ & $+0.01$ & $\checkmark$ \\
\texttt{gb\_tex\_vd1} & RQA $0.5u$ & jitter $2u$ & +0.015 [+0.009, +0.021] & \textbf{+0.016 [+0.010, +0.021]} & $+0.017$ to $+0.019$ & $+0.017$ to $+0.018$ (3/3) & \textbf{+0.014 [+0.009, +0.020]} & $+0.03$ & $+0.03$ & $\checkmark$ \\
\texttt{gb\_tex\_vd1} & RQA $1u$ & order & +0.006 [+0.001, +0.012] & \textbf{+0.007 [+0.001, +0.012]} & $+0.006$ to $+0.007$ & $+0.006$ to $+0.006$ (3/3) & \textbf{+0.007 [+0.001, +0.013]} & $+0.01$ & $+0.01$ & $\checkmark$ \\
\texttt{gb\_tex\_vd1} & RQA $1u$ & jitter $2u$ & +0.006 [+0.001, +0.011] & \textbf{+0.007 [+0.002, +0.011]} & $+0.006$ to $+0.007$ & $+0.006$ to $+0.006$ (3/3) & \textbf{+0.007 [+0.001, +0.012]} & $+0.01$ & $+0.01$ & $\checkmark$ \\
\texttt{gb\_vd1\_blg} & Classic & duration marginal & +0.033 [+0.020, +0.046] & \textbf{+0.034 [+0.022, +0.047]} & $+0.030$ to $+0.036$ & $+0.033$ to $+0.033$ (3/3) & \textbf{+0.037 [+0.009, +0.065]} & $+0.09$ & $+0.09$ & $\checkmark$ \\
\texttt{gb\_vd1\_blg} & Classic & order & $-$0.114 [$-$0.137, $-$0.091] & \textbf{$-$0.112 [$-$0.133, $-$0.091]} & $-0.114$ to $-0.110$ & $-0.114$ to $-0.114$ (3/3) & \textbf{$-$0.118 [$-$0.148, $-$0.088]} & $-0.30$ & $-0.29$ & $\checkmark$ \\
\texttt{gb\_vd1\_blg} & Classic & jitter $0.25u$ & $-$0.003 [$-$0.004, $-$0.001] & \textbf{$-$0.002 [$-$0.004, $-$0.001]} & $-0.003$ to $-0.002$ & $-0.003$ to $-0.003$ (3/3) & +0.001 [$-$0.008, +0.009] & $-0.01$ & $-0.01$ & $\checkmark$ \\
\texttt{gb\_vd1\_blg} & Classic & jitter $0.5u$ & $-$0.007 [$-$0.010, $-$0.005] & \textbf{$-$0.007 [$-$0.009, $-$0.005]} & $-0.007$ to $-0.006$ & $-0.007$ to $-0.007$ (3/3) & $-$0.003 [$-$0.012, +0.006] & $-0.02$ & $-0.02$ & $\checkmark$ \\
\texttt{gb\_vd1\_blg} & Classic & jitter $1u$ & $-$0.014 [$-$0.021, $-$0.007] & \textbf{$-$0.014 [$-$0.020, $-$0.008]} & $-0.015$ to $-0.013$ & $-0.014$ to $-0.014$ (3/3) & $-$0.009 [$-$0.021, +0.003] & $-0.04$ & $-0.04$ & $\checkmark$ \\
\texttt{gb\_vd1\_blg} & Classic & jitter $2u$ & $-$0.014 [$-$0.023, $-$0.005] & \textbf{$-$0.013 [$-$0.021, $-$0.006]} & $-0.015$ to $-0.013$ & $-0.016$ to $-0.014$ (3/3) & $-$0.005 [$-$0.023, +0.012] & $-0.04$ & $-0.03$ & $\checkmark$ \\
\texttt{gb\_vd1\_blg} & Directed & order & +0.005 [+0.001, +0.008] & \textbf{+0.005 [+0.002, +0.008]} & $+0.004$ to $+0.005$ & $+0.005$ to $+0.005$ (3/3) & \textbf{+0.006 [+0.002, +0.010]} & $+0.01$ & $+0.01$ & $\checkmark$ \\
\texttt{gb\_vd1\_blg} & Directed & jitter $2u$ & +0.002 [+0.000, +0.004] & \textbf{+0.002 [+0.001, +0.004]} & $+0.002$ to $+0.002$ & $+0.002$ to $+0.002$ (2/3) & +0.003 [$-$0.000, +0.006] & $+0.00$ & $+0.00$ & -- \\
\texttt{gb\_vd1\_blg} & Dur.\ trend & duration marginal & +0.252 [+0.199, +0.304] & \textbf{+0.257 [+0.206, +0.308]} & $+0.251$ to $+0.265$ & $+0.243$ to $+0.259$ (3/3) & \textbf{+0.226 [+0.161, +0.291]} & $+1.01$ & $+1.00$ & $\checkmark$ \\
\texttt{gb\_vd1\_blg} & FAG lattice & duration marginal & $-$0.008 [$-$0.013, $-$0.004] & \textbf{$-$0.009 [$-$0.013, $-$0.005]} & $-0.009$ to $-0.008$ & $-0.009$ to $-0.008$ (3/3) & $-$0.004 [$-$0.009, +0.002] & $-0.02$ & $-0.02$ & $\checkmark$ \\
\texttt{gb\_vd1\_blg} & FAG lattice & jitter $1u$ & +0.019 [+0.003, +0.035] & \textbf{+0.020 [+0.004, +0.035]} & $+0.017$ to $+0.021$ & $+0.017$ to $+0.019$ (3/3) & +0.011 [$-$0.000, +0.022] & $+0.04$ & $+0.04$ & $\checkmark$ \\
\texttt{gb\_vd1\_blg} & FAG lattice & jitter $2u$ & +0.027 [+0.010, +0.044] & \textbf{+0.027 [+0.011, +0.043]} & $+0.028$ to $+0.030$ & $+0.026$ to $+0.027$ (3/3) & \textbf{+0.024 [+0.009, +0.039]} & $+0.06$ & $+0.06$ & $\checkmark$ \\
\texttt{gb\_vd1\_blg} & FAG lattice & rotation & +0.029 [+0.016, +0.042] & \textbf{+0.029 [+0.018, +0.041]} & $+0.024$ to $+0.028$ & $+0.024$ to $+0.028$ (3/3) & \textbf{+0.023 [+0.009, +0.036]} & $+0.06$ & $+0.06$ & $\checkmark$ \\
\texttt{gb\_vd1\_blg} & RQA $0.5u$ & order & +0.013 [+0.006, +0.020] & \textbf{+0.013 [+0.007, +0.019]} & $+0.012$ to $+0.014$ & $+0.013$ to $+0.013$ (3/3) & \textbf{+0.014 [+0.005, +0.023]} & $+0.03$ & $+0.03$ & $\checkmark$ \\
\texttt{gb\_vd1\_blg} & RQA $0.5u$ & jitter $2u$ & +0.003 [+0.001, +0.006] & \textbf{+0.003 [+0.001, +0.006]} & $+0.005$ to $+0.006$ & $+0.004$ to $+0.004$ (3/3) & \textbf{+0.006 [+0.001, +0.010]} & $+0.01$ & $+0.01$ & $\checkmark$ \\
\texttt{gb\_vd1\_blg} & RQA $1u$ & order & +0.001 [+0.000, +0.002] & \textbf{+0.001 [+0.000, +0.002]} & $+0.001$ to $+0.001$ & $+0.001$ to $+0.001$ (3/3) & +0.002 [$-$0.001, +0.005] & $+0.00$ & $+0.00$ & $\checkmark$ \\
\texttt{gb\_vd1\_vd2} & Classic & duration marginal & +0.124 [+0.099, +0.150] & \textbf{+0.124 [+0.103, +0.144]} & $+0.124$ to $+0.135$ & $+0.122$ to $+0.127$ (3/3) & \textbf{+0.130 [+0.073, +0.187]} & $+0.88$ & $+0.90$ & $\checkmark$ \\
\texttt{gb\_vd1\_vd2} & Classic & jitter $0.5u$ & +0.004 [+0.000, +0.007] & +0.002 [$-$0.000, +0.004] & $+0.001$ to $+0.003$ & $+0.003$ to $+0.003$ (2/3) & +0.002 [$-$0.026, +0.031] & $+0.03$ & $+0.01$ & -- \\
\texttt{gb\_vd1\_vd2} & Directed & order & +0.040 [+0.015, +0.065] & \textbf{+0.042 [+0.018, +0.065]} & $+0.038$ to $+0.045$ & $+0.037$ to $+0.039$ (3/3) & \textbf{+0.059 [+0.007, +0.111]} & $+0.20$ & $+0.21$ & $\checkmark$ \\
\texttt{gb\_vd1\_vd2} & Directed & jitter $1u$ & +0.042 [+0.009, +0.075] & \textbf{+0.044 [+0.012, +0.076]} & $+0.040$ to $+0.049$ & $+0.040$ to $+0.042$ (3/3) & +0.048 [$-$0.000, +0.096] & $+0.21$ & $+0.22$ & $\checkmark$ \\
\texttt{gb\_vd1\_vd2} & Directed & jitter $2u$ & +0.069 [+0.044, +0.095] & \textbf{+0.071 [+0.045, +0.097]} & $+0.065$ to $+0.072$ & $+0.067$ to $+0.068$ (3/3) & \textbf{+0.086 [+0.040, +0.133]} & $+0.35$ & $+0.35$ & $\checkmark$ \\
\texttt{gb\_vd1\_vd2} & Dispersion & jitter $0.5u$ & +0.012 [+0.000, +0.023] & \textbf{+0.011 [+0.002, +0.021]} & $+0.010$ to $+0.013$ & $+0.012$ to $+0.013$ (3/3) & +0.035 [$-$0.015, +0.086] & $+0.06$ & $+0.06$ & $\checkmark$ \\
\texttt{gb\_vd1\_vd2} & Dispersion & jitter $1u$ & +0.045 [+0.026, +0.063] & \textbf{+0.043 [+0.027, +0.060]} & $+0.039$ to $+0.043$ & $+0.040$ to $+0.048$ (3/3) & \textbf{+0.069 [+0.005, +0.134]} & $+0.21$ & $+0.21$ & $\checkmark$ \\
\texttt{gb\_vd1\_vd2} & Dispersion & jitter $2u$ & +0.116 [+0.085, +0.148] & \textbf{+0.113 [+0.082, +0.143]} & $+0.117$ to $+0.123$ & $+0.113$ to $+0.120$ (3/3) & \textbf{+0.102 [+0.038, +0.166]} & $+0.56$ & $+0.55$ & $\checkmark$ \\
\texttt{gb\_vd1\_vd2} & Dispersion & rotation & +0.098 [+0.068, +0.128] & \textbf{+0.096 [+0.069, +0.124]} & $+0.108$ to $+0.114$ & $+0.100$ to $+0.103$ (3/3) & \textbf{+0.108 [+0.045, +0.170]} & $+0.47$ & $+0.47$ & $\checkmark$ \\
\texttt{gb\_vd1\_vd2} & Dur.\ trend & duration coupling & +0.158 [+0.131, +0.185] & \textbf{+0.162 [+0.140, +0.184]} & $+0.154$ to $+0.162$ & $+0.157$ to $+0.158$ (3/3) & \textbf{+0.176 [+0.133, +0.219]} & $+0.54$ & $+0.55$ & $\checkmark$ \\
\texttt{gb\_vd1\_vd2} & Dur.\ trend & duration marginal & +0.286 [+0.227, +0.345] & \textbf{+0.292 [+0.233, +0.351]} & $+0.280$ to $+0.294$ & $+0.289$ to $+0.297$ (3/3) & \textbf{+0.272 [+0.216, +0.329]} & $+0.98$ & $+0.98$ & $\checkmark$ \\
\texttt{gb\_vd1\_vd2} & Dur.\ trend & order & +0.157 [+0.127, +0.187] & \textbf{+0.161 [+0.134, +0.188]} & $+0.150$ to $+0.160$ & $+0.157$ to $+0.157$ (3/3) & \textbf{+0.178 [+0.133, +0.224]} & $+0.54$ & $+0.54$ & $\checkmark$ \\
\texttt{gb\_vd1\_vd2} & FAG lattice & duration coupling & $-$0.015 [$-$0.029, $-$0.001] & $-$0.012 [$-$0.028, +0.003] & $-0.017$ to $-0.012$ & $-0.017$ to $-0.015$ (1/3) & $-$0.016 [$-$0.046, +0.013] & $-0.09$ & $-0.07$ & -- \\
\texttt{gb\_vd1\_vd2} & FAG lattice & duration marginal & $-$0.020 [$-$0.031, $-$0.009] & \textbf{$-$0.019 [$-$0.032, $-$0.007]} & $-0.027$ to $-0.020$ & $-0.023$ to $-0.020$ (3/3) & $-$0.014 [$-$0.038, +0.011] & $-0.12$ & $-0.11$ & $\checkmark$ \\
\texttt{gb\_vd1\_vd2} & FAG lattice & jitter $2u$ & +0.055 [+0.013, +0.096] & \textbf{+0.058 [+0.017, +0.099]} & $+0.045$ to $+0.053$ & $+0.046$ to $+0.052$ (3/3) & \textbf{+0.100 [+0.043, +0.156]} & $+0.33$ & $+0.34$ & $\checkmark$ \\
\texttt{gb\_vd1\_vd2} & RQA $0.5u$ & order & +0.035 [+0.011, +0.060] & \textbf{+0.034 [+0.013, +0.054]} & $+0.032$ to $+0.040$ & $+0.033$ to $+0.035$ (3/3) & +0.044 [$-$0.004, +0.092] & $+0.20$ & $+0.19$ & $\checkmark$ \\
\texttt{gb\_vd1\_vd2} & RQA $0.5u$ & jitter $1u$ & +0.046 [+0.003, +0.090] & \textbf{+0.044 [+0.000, +0.087]} & $+0.033$ to $+0.040$ & $+0.035$ to $+0.039$ (3/3) & +0.049 [$-$0.023, +0.121] & $+0.26$ & $+0.25$ & $\checkmark$ \\
\texttt{gb\_vd1\_vd2} & RQA $0.5u$ & jitter $2u$ & +0.108 [+0.071, +0.145] & \textbf{+0.106 [+0.071, +0.140]} & $+0.107$ to $+0.114$ & $+0.106$ to $+0.115$ (3/3) & \textbf{+0.106 [+0.041, +0.170]} & $+0.61$ & $+0.61$ & $\checkmark$ \\
\texttt{gb\_vd1\_vd2} & RQA $1u$ & order & +0.073 [+0.043, +0.103] & \textbf{+0.073 [+0.045, +0.101]} & $+0.069$ to $+0.076$ & $+0.072$ to $+0.073$ (3/3) & \textbf{+0.086 [+0.032, +0.140]} & $+0.33$ & $+0.33$ & $\checkmark$ \\
\texttt{gb\_vd1\_vd2} & RQA $1u$ & jitter $1u$ & +0.043 [+0.007, +0.078] & \textbf{+0.041 [+0.007, +0.075]} & $+0.041$ to $+0.048$ & $+0.040$ to $+0.042$ (3/3) & +0.044 [$-$0.005, +0.093] & $+0.19$ & $+0.19$ & $\checkmark$ \\
\texttt{gb\_vd1\_vd2} & RQA $1u$ & jitter $2u$ & +0.118 [+0.080, +0.156] & \textbf{+0.115 [+0.077, +0.152]} & $+0.108$ to $+0.115$ & $+0.110$ to $+0.119$ (3/3) & \textbf{+0.125 [+0.054, +0.195]} & $+0.54$ & $+0.52$ & $\checkmark$ \\
\texttt{gb\_vd1\_halves} & Classic & order & $-$0.090 [$-$0.110, $-$0.069] & \textbf{$-$0.091 [$-$0.111, $-$0.071]} & $-0.098$ to $-0.088$ & $-0.089$ to $-0.087$ (3/3) & $-$0.055 [$-$0.114, +0.003] & $-1.40$ & $-1.45$ & $\checkmark$ \\
\texttt{gb\_vd1\_halves} & Directed & order & +0.036 [+0.017, +0.054] & \textbf{+0.036 [+0.019, +0.052]} & $+0.032$ to $+0.036$ & $+0.037$ to $+0.037$ (3/3) & +0.027 [$-$0.003, +0.058] & $+0.12$ & $+0.12$ & $\checkmark$ \\
\texttt{gb\_vd1\_halves} & Directed & jitter $0.5u$ & +0.033 [+0.014, +0.052] & \textbf{+0.031 [+0.011, +0.052]} & $+0.026$ to $+0.031$ & $+0.031$ to $+0.033$ (3/3) & +0.029 [$-$0.007, +0.065] & $+0.11$ & $+0.10$ & $\checkmark$ \\
\texttt{gb\_vd1\_halves} & Directed & jitter $1u$ & +0.090 [+0.060, +0.121] & \textbf{+0.089 [+0.061, +0.117]} & $+0.087$ to $+0.093$ & $+0.087$ to $+0.093$ (3/3) & \textbf{+0.102 [+0.060, +0.145]} & $+0.29$ & $+0.29$ & $\checkmark$ \\
\texttt{gb\_vd1\_halves} & Directed & jitter $2u$ & +0.191 [+0.161, +0.220] & \textbf{+0.190 [+0.162, +0.218]} & $+0.167$ to $+0.172$ & $+0.186$ to $+0.187$ (3/3) & \textbf{+0.192 [+0.140, +0.243]} & $+0.62$ & $+0.62$ & $\checkmark$ \\
\texttt{gb\_vd1\_halves} & Dispersion & jitter $0.5u$ & +0.007 [+0.002, +0.012] & \textbf{+0.007 [+0.002, +0.011]} & $+0.006$ to $+0.008$ & $+0.007$ to $+0.007$ (3/3) & +0.001 [$-$0.011, +0.013] & $+0.02$ & $+0.02$ & $\checkmark$ \\
\texttt{gb\_vd1\_halves} & Dispersion & jitter $1u$ & +0.029 [+0.018, +0.041] & \textbf{+0.029 [+0.020, +0.038]} & $+0.031$ to $+0.033$ & $+0.030$ to $+0.032$ (3/3) & \textbf{+0.028 [+0.010, +0.046]} & $+0.07$ & $+0.07$ & $\checkmark$ \\
\texttt{gb\_vd1\_halves} & Dispersion & jitter $2u$ & +0.130 [+0.111, +0.150] & \textbf{+0.129 [+0.110, +0.149]} & $+0.123$ to $+0.126$ & $+0.122$ to $+0.129$ (3/3) & \textbf{+0.123 [+0.097, +0.150]} & $+0.31$ & $+0.31$ & $\checkmark$ \\
\texttt{gb\_vd1\_halves} & Dispersion & rotation & +0.249 [+0.224, +0.275] & \textbf{+0.247 [+0.223, +0.271]} & $+0.249$ to $+0.253$ & $+0.246$ to $+0.249$ (3/3) & \textbf{+0.177 [+0.139, +0.215]} & $+0.59$ & $+0.58$ & $\checkmark$ \\
\texttt{gb\_vd1\_halves} & Dur.\ trend & duration coupling & +0.085 [+0.040, +0.130] & \textbf{+0.086 [+0.042, +0.129]} & $+0.078$ to $+0.090$ & $+0.078$ to $+0.082$ (3/3) & +0.060 [$-$0.002, +0.121] & $+0.96$ & $+0.91$ & $\checkmark$ \\
\texttt{gb\_vd1\_halves} & Dur.\ trend & duration marginal & +0.082 [+0.044, +0.120] & \textbf{+0.088 [+0.051, +0.125]} & $+0.077$ to $+0.085$ & $+0.089$ to $+0.095$ (3/3) & +0.024 [$-$0.051, +0.098] & $+0.93$ & $+0.94$ & $\checkmark$ \\
\texttt{gb\_vd1\_halves} & Dur.\ trend & order & +0.076 [+0.030, +0.122] & \textbf{+0.079 [+0.031, +0.126]} & $+0.072$ to $+0.080$ & $+0.080$ to $+0.081$ (3/3) & +0.033 [$-$0.045, +0.112] & $+0.86$ & $+0.84$ & $\checkmark$ \\
\texttt{gb\_vd1\_halves} & FAG lattice & jitter $2u$ & +0.088 [+0.053, +0.123] & \textbf{+0.087 [+0.052, +0.121]} & $+0.084$ to $+0.090$ & $+0.082$ to $+0.082$ (3/3) & \textbf{+0.111 [+0.049, +0.173]} & $+0.37$ & $+0.36$ & $\checkmark$ \\
\texttt{gb\_vd1\_halves} & RQA $0.5u$ & order & +0.046 [+0.023, +0.069] & \textbf{+0.046 [+0.025, +0.068]} & $+0.043$ to $+0.049$ & $+0.047$ to $+0.048$ (3/3) & \textbf{+0.055 [+0.014, +0.096]} & $+0.19$ & $+0.19$ & $\checkmark$ \\
\texttt{gb\_vd1\_halves} & RQA $0.5u$ & jitter $1u$ & +0.079 [+0.051, +0.106] & \textbf{+0.079 [+0.051, +0.107]} & $+0.068$ to $+0.074$ & $+0.070$ to $+0.071$ (3/3) & \textbf{+0.080 [+0.031, +0.129]} & $+0.32$ & $+0.33$ & $\checkmark$ \\
\texttt{gb\_vd1\_halves} & RQA $0.5u$ & jitter $2u$ & +0.143 [+0.097, +0.189] & \textbf{+0.143 [+0.099, +0.188]} & $+0.137$ to $+0.142$ & $+0.145$ to $+0.150$ (3/3) & \textbf{+0.157 [+0.101, +0.214]} & $+0.59$ & $+0.59$ & $\checkmark$ \\
\texttt{gb\_vd1\_halves} & RQA $1u$ & order & +0.088 [+0.063, +0.112] & \textbf{+0.089 [+0.066, +0.111]} & $+0.084$ to $+0.090$ & $+0.085$ to $+0.086$ (3/3) & \textbf{+0.113 [+0.077, +0.149]} & $+0.27$ & $+0.27$ & $\checkmark$ \\
\texttt{gb\_vd1\_halves} & RQA $1u$ & jitter $0.5u$ & +0.026 [+0.006, +0.046] & \textbf{+0.026 [+0.007, +0.045]} & $+0.018$ to $+0.021$ & $+0.028$ to $+0.031$ (3/3) & +0.025 [$-$0.011, +0.061] & $+0.08$ & $+0.08$ & $\checkmark$ \\
\texttt{gb\_vd1\_halves} & RQA $1u$ & jitter $1u$ & +0.088 [+0.054, +0.121] & \textbf{+0.089 [+0.059, +0.118]} & $+0.086$ to $+0.090$ & $+0.089$ to $+0.092$ (3/3) & \textbf{+0.106 [+0.071, +0.142]} & $+0.27$ & $+0.27$ & $\checkmark$ \\
\texttt{gb\_vd1\_halves} & RQA $1u$ & jitter $2u$ & +0.182 [+0.151, +0.214] & \textbf{+0.183 [+0.154, +0.213]} & $+0.170$ to $+0.174$ & $+0.179$ to $+0.183$ (3/3) & \textbf{+0.212 [+0.172, +0.253]} & $+0.56$ & $+0.57$ & $\checkmark$ \\
\texttt{gb\_tex\_s1s2} & RQA $0.5u$ & order & +0.093 [+0.013, +0.174] & \textbf{+0.095 [+0.017, +0.173]} & $+0.077$ to $+0.097$ & $+0.064$ to $+0.083$ (1/3) & +0.066 [$-$0.037, +0.168] & $+1.23$ & $+1.19$ & -- \\
\texttt{gb\_tex\_s1s2} & RQA $0.5u$ & jitter $0.5u$ & +0.071 [+0.020, +0.123] & \textbf{+0.071 [+0.017, +0.125]} & $+0.048$ to $+0.062$ & $+0.057$ to $+0.065$ (0/3) & +0.035 [$-$0.065, +0.135] & $+0.94$ & $+0.89$ & -- \\
\texttt{sb\_native} & RQA $0.5u$ & jitter $2u$ & +0.088 [+0.012, +0.164] & \textbf{+0.096 [+0.009, +0.182]} & $+0.104$ to $+0.125$ & $+0.103$ to $+0.117$ (3/3) & $-$0.002 [$-$0.084, +0.081] & $+0.79$ & $+0.77$ & $\checkmark$ \\
\texttt{sb\_native} & RQA $0.5u$ & jitter $1u$ & +0.079 [+0.007, +0.150] & +0.085 [$-$0.005, +0.174] & $+0.055$ to $+0.072$ & $+0.069$ to $+0.076$ (2/3) & $-$0.020 [$-$0.102, +0.062] & $+0.71$ & $+0.68$ & -- \\
\texttt{sb\_native} & RQA $1u$ & jitter $2u$ & +0.140 [+0.045, +0.235] & \textbf{+0.141 [+0.035, +0.246]} & $+0.139$ to $+0.171$ & $+0.154$ to $+0.155$ (3/3) & \textbf{+0.100 [+0.009, +0.190]} & $+0.85$ & $+0.81$ & $\checkmark$ \\
\texttt{sb\_native} & RQA $1u$ & jitter $1u$ & +0.086 [+0.016, +0.157] & \textbf{+0.085 [+0.009, +0.161]} & $+0.068$ to $+0.087$ & $+0.080$ to $+0.084$ (3/3) & +0.071 [$-$0.003, +0.145] & $+0.52$ & $+0.49$ & $\checkmark$ \\
\texttt{em\_program} & Dispersion & rotation & +0.139 [+0.078, +0.199] & \textbf{+0.160 [+0.108, +0.213]} & $+0.137$ to $+0.165$ & $+0.134$ to $+0.142$ (3/3) & \textbf{+0.166 [+0.088, +0.244]} & $+0.30$ & $+0.34$ & $\checkmark$ \\
\texttt{em\_halves} & Classic & order & +0.172 [+0.112, +0.231] & \textbf{+0.171 [+0.120, +0.222]} & $+0.164$ to $+0.191$ & $+0.171$ to $+0.174$ (3/3) & \textbf{+0.128 [+0.033, +0.222]} & $+0.62$ & $+0.62$ & $\checkmark$ \\
\texttt{em\_halves} & FAG AOI & jitter $2u$ & +0.113 [+0.014, +0.212] & \textbf{+0.111 [+0.023, +0.198]} & $+0.101$ to $+0.122$ & $+0.111$ to $+0.115$ (3/3) & +0.115 [$-$0.007, +0.237] & $+0.62$ & $+0.61$ & -- \\
\texttt{em\_halves} & FAG AOI & rotation & +0.098 [+0.031, +0.166] & \textbf{+0.093 [+0.028, +0.158]} & $+0.081$ to $+0.100$ & $+0.102$ to $+0.114$ (3/3) & +0.058 [$-$0.060, +0.176] & $+0.54$ & $+0.51$ & -- \\
\texttt{em\_halves} & RQA $0.5u$ & jitter $1u$ & $-$0.085 [$-$0.152, $-$0.018] & \textbf{$-$0.085 [$-$0.145, $-$0.025]} & $-0.089$ to $-0.073$ & $-0.085$ to $-0.078$ (3/3) & $-$0.089 [$-$0.189, +0.011] & $-0.77$ & $-0.75$ & $\checkmark$ \\
\texttt{em\_halves} & RQA $1u$ & jitter $1u$ & $-$0.052 [$-$0.098, $-$0.006] & \textbf{$-$0.052 [$-$0.086, $-$0.017]} & $-0.063$ to $-0.045$ & $-0.050$ to $-0.045$ (2/3) & $-$0.053 [$-$0.127, +0.020] & $-0.29$ & $-0.28$ & -- \\
\end{longtable}

}

\begin{table}[!htbp]\centering\scriptsize
\caption{Bootstrap standard error against the empirical spread of the estimate in the calibration campaign ($R = 6$, $B = 60$; 40 datasets per design and size). SE ratio: mean bootstrap SE divided by the SD of the estimate across datasets, per contrast; median and range over the contrasts of the cell. Coverage of the target by normal intervals at nominal 80\%, 90\% and 95\%.}
\begin{tabular}{llrrrrr}
\toprule
Design & Size & Contrasts & SE ratio, median (range) & Coverage at 80\% & at 90\% & at 95\% \\
\midrule
no signal & 65 $\times$ 2 (B) & 8 & 1.66 (1.35--2.04) & 0.959 & 0.988 & 0.997 \\
no signal & 95 $\times$ 20 (B) & 8 & 1.48 (1.23--2.03) & 0.922 & 0.978 & 0.988 \\
no signal & 112 $\times$ 2 (W) & 8 & 1.80 (1.66--2.02) & 0.963 & 0.994 & 1.000 \\
no signal & 60 $\times$ 4 (W) & 8 & 1.80 (1.48--2.22) & 0.966 & 0.987 & 1.000 \\
signal elsewhere & 65 $\times$ 2 (B) & 5 & 2.40 (1.94--3.10) & 0.985 & 0.990 & 1.000 \\
signal elsewhere & 95 $\times$ 20 (B) & 5 & 1.46 (1.28--1.65) & 0.925 & 0.985 & 1.000 \\
signal elsewhere & 112 $\times$ 2 (W) & 5 & 2.35 (2.04--2.82) & 0.990 & 1.000 & 1.000 \\
signal elsewhere & 60 $\times$ 4 (W) & 5 & 2.65 (1.76--3.37) & 0.995 & 1.000 & 1.000 \\
signal elsewhere, no offset & 65 $\times$ 2 (B) & 3 & 2.40 (1.87--2.58) & 1.000 & 1.000 & 1.000 \\
signal elsewhere, no offset & 95 $\times$ 20 (B) & 3 & 1.29 (1.28--1.66) & 0.925 & 0.967 & 0.983 \\
signal elsewhere, no offset & 112 $\times$ 2 (W) & 3 & 1.97 (1.57--2.82) & 0.975 & 1.000 & 1.000 \\
signal elsewhere, no offset & 60 $\times$ 4 (W) & 3 & 2.19 (1.67--3.37) & 0.975 & 1.000 & 1.000 \\
\bottomrule
\end{tabular}

\end{table}

\begin{table}[!htbp]\centering\scriptsize
\caption{Corrected refit bootstrap (Equation 1) against the direct refit bootstrap, which averages all $R = 6$ transformed replicates inside each of the $B = 60$ draws, on the same 200 datasets per cell. Coverage of the target, mean width, and the share of intervals excluding zero.}
\begin{tabular}{llllrrrrrr}
\toprule
 & & & & \multicolumn{2}{c}{Coverage} & \multicolumn{2}{c}{Width} & \multicolumn{2}{c}{Excl.\ 0} \\
\cmidrule(lr){5-6}\cmidrule(lr){7-8}\cmidrule(lr){9-10}
Design & Size & Descr. & Contrast & direct & corr. & direct & corr. & direct & corr. \\
\midrule
no signal & 65 $\times$ 2 (B) & rqa & order & 0.985 & 0.990 & 0.367 & 0.434 & 0.015 & 0.010 \\
no signal & 65 $\times$ 2 (B) & rqa & jitter 160 px & 0.995 & 1.000 & 0.419 & 0.505 & 0.005 & 0.000 \\
no signal & 65 $\times$ 2 (B) & rqa & jitter 40 px & 0.985 & 0.990 & 0.417 & 0.498 & 0.015 & 0.010 \\
no signal & 60 $\times$ 4 (W) & classic & order & 0.995 & 0.985 & 0.241 & 0.286 & 0.005 & 0.015 \\
no signal & 60 $\times$ 4 (W) & rqa & order & 0.990 & 0.995 & 0.272 & 0.324 & 0.010 & 0.005 \\
no signal & 60 $\times$ 4 (W) & rqa & jitter 160 px & 0.995 & 1.000 & 0.313 & 0.379 & 0.005 & 0.005 \\
no signal & 60 $\times$ 4 (W) & rqa & jitter 40 px & 0.990 & 1.000 & 0.314 & 0.376 & 0.010 & 0.000 \\
signal elsewhere & 65 $\times$ 2 (B) & rqa & order & 1.000 & 1.000 & 0.125 & 0.161 & 0.000 & 0.000 \\
signal elsewhere & 65 $\times$ 2 (B) & rqa & jitter 40 px & 1.000 & 1.000 & 0.219 & 0.345 & 1.000 & 0.955 \\
signal elsewhere & 60 $\times$ 4 (W) & rqa & order & 1.000 & 1.000 & 0.070 & 0.089 & 0.000 & 0.000 \\
signal elsewhere & 60 $\times$ 4 (W) & rqa & jitter 40 px & 0.995 & 1.000 & 0.147 & 0.238 & 1.000 & 1.000 \\
signal elsewhere, no offset & 60 $\times$ 4 (W) & rqa & order & 1.000 & 1.000 & 0.070 & 0.089 & 0.000 & 0.000 \\
\bottomrule
\end{tabular}

\end{table}

{\scriptsize\setlength{\LTcapwidth}{\linewidth}
\begin{longtable}{lllrrrrr}
\caption{Extended simulation designs ($R = 6$, $B = 60$; 40 datasets per row). Target: the reference value (60 datasets with disjoint seeds, $R = 20$) with its Monte Carlo standard error, or zero by exchangeability. Unequal: one to four trials per class and participant (W) or per participant (B); 1:3 and 16:49: class imbalance; heterogeneous: participant offset SD 100 px and participant-specific scaling.}\\
\toprule
Design & Size & Contrast & Target $\theta$ (SE) & Coverage & Excl.\ 0 & Width & Obs.\ AUC \\
\midrule
\endfirsthead
\toprule
Design & Size & Contrast & Target $\theta$ (SE) & Coverage & Excl.\ 0 & Width & Obs.\ AUC \\
\midrule
\endhead
\bottomrule
\endfoot
order hides spread (classic) & 65 $\times$ 2 (B) & duration marginal & +0.006 (0.009) & 0.975 & 0.025 & 0.397 & 0.520 \\
order hides spread (classic) & 65 $\times$ 2 (B) & order & $-$0.499 (0.010) & 1.000 & 1.000 & 0.405 & 0.520 \\
order hides spread (classic) & 60 $\times$ 4 (W) & duration marginal & $-$0.001 (0.005) & 1.000 & 0.000 & 0.296 & 0.509 \\
order hides spread (classic) & 60 $\times$ 4 (W) & order & $-$0.497 (0.008) & 1.000 & 1.000 & 0.292 & 0.509 \\
no signal (rqa) & 65 $\times$ 2, heterog.\ (B) & order & 0 (by arg.) & 1.000 & 0.000 & 0.429 & 0.493 \\
no signal (rqa) & 65 $\times$ 2, heterog.\ (B) & jitter 160 px & +0.001 (0.012) & 1.000 & 0.000 & 0.498 & 0.493 \\
no signal (rqa) & 65 $\times$ 2, heterog.\ (B) & jitter 40 px & +0.005 (0.011) & 1.000 & 0.000 & 0.451 & 0.493 \\
no signal (rqa) & 65, 16:49 (B) & order & 0 (by arg.) & 0.925 & 0.075 & 0.502 & 0.474 \\
no signal (rqa) & 65, 16:49 (B) & jitter 160 px & +0.005 (0.012) & 1.000 & 0.000 & 0.605 & 0.474 \\
no signal (rqa) & 65, 16:49 (B) & jitter 40 px & +0.005 (0.012) & 1.000 & 0.000 & 0.573 & 0.474 \\
no signal (rqa) & 65, unequal (B) & order & 0 (by arg.) & 1.000 & 0.000 & 0.389 & 0.508 \\
no signal (rqa) & 65, unequal (B) & jitter 160 px & +0.014 (0.008) & 1.000 & 0.000 & 0.455 & 0.508 \\
no signal (rqa) & 65, unequal (B) & jitter 40 px & +0.016 (0.008) & 1.000 & 0.000 & 0.446 & 0.508 \\
no signal (rqa) & 60 $\times$ 4, heterog.\ (W) & order & 0 (by arg.) & 1.000 & 0.000 & 0.377 & 0.508 \\
no signal (rqa) & 60 $\times$ 4, heterog.\ (W) & jitter 160 px & $-$0.009 (0.006) & 1.000 & 0.000 & 0.375 & 0.508 \\
no signal (rqa) & 60 $\times$ 4, heterog.\ (W) & jitter 40 px & $-$0.007 (0.006) & 1.000 & 0.000 & 0.362 & 0.508 \\
no signal (rqa) & 60, 1:3 (W) & order & 0 (by arg.) & 1.000 & 0.000 & 0.393 & 0.510 \\
no signal (rqa) & 60, 1:3 (W) & jitter 160 px & +0.010 (0.009) & 1.000 & 0.000 & 0.445 & 0.510 \\
no signal (rqa) & 60, 1:3 (W) & jitter 40 px & +0.007 (0.009) & 1.000 & 0.000 & 0.432 & 0.510 \\
no signal (rqa) & 60, unequal (W) & order & 0 (by arg.) & 1.000 & 0.000 & 0.295 & 0.490 \\
no signal (rqa) & 60, unequal (W) & jitter 160 px & +0.012 (0.006) & 1.000 & 0.000 & 0.349 & 0.490 \\
no signal (rqa) & 60, unequal (W) & jitter 40 px & +0.008 (0.006) & 1.000 & 0.000 & 0.334 & 0.490 \\
weak order contrast (directed) & 65 $\times$ 2 (B) & order & +0.118 (0.008) & 1.000 & 0.075 & 0.424 & 0.632 \\
weak order contrast (directed) & 65 $\times$ 2 (B) & jitter 40 px & +0.015 (0.007) & 1.000 & 0.000 & 0.425 & 0.632 \\
weak order contrast (directed) & 65 $\times$ 2 (B) & rotation & +0.000 (0.006) & 1.000 & 0.000 & 0.373 & 0.632 \\
weak order contrast (directed) & 60 $\times$ 4 (W) & order & +0.128 (0.007) & 0.975 & 0.350 & 0.307 & 0.639 \\
weak order contrast (directed) & 60 $\times$ 4 (W) & jitter 40 px & +0.015 (0.005) & 1.000 & 0.000 & 0.298 & 0.639 \\
weak order contrast (directed) & 60 $\times$ 4 (W) & rotation & $-$0.004 (0.005) & 1.000 & 0.000 & 0.265 & 0.639 \\
weak spread contrast (rqa) & 65 $\times$ 2, heterog.\ (B) & order & 0 (by arg.) & 1.000 & 0.000 & 0.282 & 0.671 \\
weak spread contrast (rqa) & 65 $\times$ 2, heterog.\ (B) & jitter 40 px & +0.127 (0.007) & 1.000 & 0.125 & 0.412 & 0.671 \\
weak spread contrast (rqa) & 65, 16:49 (B) & order & 0 (by arg.) & 1.000 & 0.000 & 0.172 & 0.871 \\
weak spread contrast (rqa) & 65, 16:49 (B) & jitter 40 px & +0.315 (0.005) & 1.000 & 0.975 & 0.394 & 0.871 \\
weak spread contrast (rqa) & 65, unequal (B) & order & 0 (by arg.) & 1.000 & 0.000 & 0.106 & 0.878 \\
weak spread contrast (rqa) & 65, unequal (B) & jitter 40 px & +0.310 (0.004) & 1.000 & 1.000 & 0.302 & 0.878 \\
weak spread contrast (rqa) & 65 $\times$ 2 (B) & order & 0 (by arg.) & 1.000 & 0.000 & 0.133 & 0.882 \\
weak spread contrast (rqa) & 65 $\times$ 2 (B) & jitter 40 px & +0.306 (0.004) & 1.000 & 1.000 & 0.334 & 0.882 \\
weak spread contrast (rqa) & 60 $\times$ 4, heterog.\ (W) & order & 0 (by arg.) & 1.000 & 0.000 & 0.152 & 0.693 \\
weak spread contrast (rqa) & 60 $\times$ 4, heterog.\ (W) & jitter 40 px & +0.143 (0.003) & 1.000 & 0.600 & 0.250 & 0.693 \\
weak spread contrast (rqa) & 60, 1:3 (W) & order & 0 (by arg.) & 1.000 & 0.000 & 0.093 & 0.872 \\
weak spread contrast (rqa) & 60, 1:3 (W) & jitter 40 px & +0.302 (0.004) & 0.975 & 1.000 & 0.290 & 0.872 \\
weak spread contrast (rqa) & 60, unequal (W) & order & 0 (by arg.) & 1.000 & 0.000 & 0.059 & 0.881 \\
weak spread contrast (rqa) & 60, unequal (W) & jitter 40 px & +0.300 (0.003) & 1.000 & 1.000 & 0.209 & 0.881 \\
weak spread contrast (rqa) & 60 $\times$ 4 (W) & order & 0 (by arg.) & 1.000 & 0.000 & 0.070 & 0.874 \\
weak spread contrast (rqa) & 60 $\times$ 4 (W) & jitter 40 px & +0.302 (0.003) & 1.000 & 1.000 & 0.247 & 0.874 \\
signal elsewhere, no offset (rqa) & 65 $\times$ 2, heterog.\ (B) & order & 0 (by arg.) & 1.000 & 0.000 & 0.299 & 0.637 \\
signal elsewhere, no offset (rqa) & 65, 16:49 (B) & order & 0 (by arg.) & 1.000 & 0.000 & 0.213 & 0.824 \\
signal elsewhere, no offset (rqa) & 65, unequal (B) & order & 0 (by arg.) & 1.000 & 0.000 & 0.126 & 0.836 \\
signal elsewhere, no offset (rqa) & 60 $\times$ 4, heterog.\ (W) & order & 0 (by arg.) & 1.000 & 0.000 & 0.177 & 0.660 \\
signal elsewhere, no offset (rqa) & 60, 1:3 (W) & order & 0 (by arg.) & 1.000 & 0.000 & 0.113 & 0.833 \\
signal elsewhere, no offset (rqa) & 60, unequal (W) & order & 0 (by arg.) & 1.000 & 0.000 & 0.072 & 0.841 \\
\end{longtable}

}

\begin{table}[!htbp]\centering\scriptsize
\caption{Controls for the EMIP program and correctness contrasts (within-fold AUC). Program contrast: AUC on real gaze and on synthetic gaze, in which every trial is replaced by the same number of fixations drawn uniformly over the bounding box of its program's code tokens. Correctness, both programs: AUC over all correct--incorrect pairs of a test fold, the mean of its permutation distribution when labels are permuted within programs, and the AUC over pairs from the same program and test fold with a refit participant-bootstrap interval (200 draws) and a one-sided permutation $p$-value (1000 permutations within programs). Descriptors with no spatial features are not affected by the synthetic control and are marked with a dash.}
\begin{tabular}{lrrrrr}
\toprule
 & \multicolumn{2}{c}{Program contrast} & \multicolumn{3}{c}{Correctness, both programs} \\
\cmidrule(lr){2-3}\cmidrule(lr){4-6}
Descriptor & Real gaze & Synthetic gaze & All pairs & Perm.\ mean & Same-program pairs [95\% interval] \\
\midrule
Count & -- & -- & 0.495 & 0.515 & 0.517 [0.317, 0.646] ($p = 0.455$) \\
Dispersion & 0.957 & 1.000 & 0.716 & 0.710 & 0.490 [0.398, 0.581] ($p = 0.594$) \\
Classic & -- & -- & 0.501 & 0.505 & 0.523 [0.319, 0.625] ($p = 0.394$) \\
Dur.\ trend & -- & -- & 0.511 & 0.502 & 0.538 [0.373, 0.640] ($p = 0.262$) \\
RQA $0.5u$ & 0.580 & 0.604 & 0.571 & 0.517 & 0.527 [0.403, 0.645] ($p = 0.304$) \\
RQA $1u$ & 0.547 & 0.648 & 0.525 & 0.502 & 0.509 [0.384, 0.636] ($p = 0.462$) \\
Directed & 0.655 & 0.554 & 0.580 & 0.540 & 0.562 [0.410, 0.629] ($p = 0.130$) \\
FAG lattice & 0.658 & 0.859 & 0.664 & 0.566 & 0.577 [0.484, 0.685] ($p = 0.042$) \\
FAG AOI & 0.778 & 0.879 & 0.635 & 0.619 & 0.559 [0.441, 0.666] ($p = 0.087$) \\
FAG AOI pad. & 0.729 & 0.789 & 0.571 & 0.591 & 0.494 [0.345, 0.603] ($p = 0.575$) \\
\bottomrule
\end{tabular}

\end{table}

\begin{table}[!htbp]\centering\scriptsize
\caption{EMIP expertise: AUC of each of the ten test folds and of the corresponding training set, for three descriptors (linear SVM; class-balanced folds).}
\begin{tabular}{lrrrrrr}
\toprule
Fold & \multicolumn{2}{c}{Classic} & \multicolumn{2}{c}{Directed} & \multicolumn{2}{c}{FAG AOI} \\
 & test & train & test & train & test & train \\
\midrule
1 & 0.686 & 0.532 & 0.400 & 0.593 & 0.229 & 0.588 \\
2 & 0.167 & 0.582 & 0.033 & 0.616 & 0.300 & 0.533 \\
3 & 0.143 & 0.595 & 0.381 & 0.589 & 0.095 & 0.598 \\
4 & 0.929 & 0.497 & 0.964 & 0.485 & 0.107 & 0.523 \\
5 & 0.429 & 0.551 & 0.257 & 0.594 & 0.200 & 0.595 \\
6 & 0.000 & 0.605 & 0.133 & 0.608 & 0.467 & 0.532 \\
7 & 0.300 & 0.562 & 0.500 & 0.576 & 0.100 & 0.549 \\
8 & 0.733 & 0.528 & 0.367 & 0.539 & 0.533 & 0.531 \\
9 & 0.214 & 0.533 & 0.214 & 0.538 & 0.286 & 0.542 \\
10 & 0.844 & 0.528 & 0.438 & 0.569 & 0.344 & 0.567 \\
\midrule
Mean & 0.444 & 0.551 & 0.369 & 0.571 & 0.266 & 0.556 \\
\bottomrule
\end{tabular}

\end{table}

\begin{table}[!htbp]\centering\scriptsize
\caption{Share of dwell time on the AOI triangulation, negative / positive class, for every EMIP problem and transformation (mean over trials of the mean over $R = 6$ replicates). Fixations off the triangulation are dropped by the AOI graph. On the lattice graph and on the padded AOI graph the share was at least 0.998 in every cell.}
\resizebox{\linewidth}{!}{\begin{tabular}{lcccccc}
\toprule
Transformation & \texttt{em\_program} & \texttt{em\_halves} & \texttt{em\_expertise} & \texttt{em\_correct} & \texttt{em\_correctrect} & \texttt{em\_correctveh} \\
\midrule
none (observed) & 0.974 / 0.975 & 0.960 / 0.990 & 0.977 / 0.973 & 0.973 / 0.974 & 0.972 / 0.972 & 0.978 / 0.975 \\
order & 0.974 / 0.975 & 0.960 / 0.990 & 0.977 / 0.973 & 0.973 / 0.974 & 0.972 / 0.972 & 0.978 / 0.975 \\
duration coupling & 0.975 / 0.977 & 0.962 / 0.990 & 0.978 / 0.975 & 0.974 / 0.975 & 0.973 / 0.972 & 0.978 / 0.977 \\
duration marginal & 0.975 / 0.977 & 0.962 / 0.990 & 0.978 / 0.975 & 0.974 / 0.975 & 0.973 / 0.972 & 0.979 / 0.977 \\
rotation & 0.901 / 0.767 & 0.788 / 0.888 & 0.838 / 0.814 & 0.864 / 0.816 & 0.897 / 0.882 & 0.742 / 0.785 \\
jitter 0.25$u$ & 0.970 / 0.970 & 0.956 / 0.983 & 0.971 / 0.968 & 0.968 / 0.970 & 0.968 / 0.968 & 0.971 / 0.970 \\
jitter 0.5$u$ & 0.959 / 0.957 & 0.945 / 0.971 & 0.956 / 0.958 & 0.955 / 0.958 & 0.955 / 0.960 & 0.957 / 0.958 \\
jitter 1$u$ & 0.927 / 0.922 & 0.911 / 0.937 & 0.920 / 0.923 & 0.921 / 0.927 & 0.920 / 0.931 & 0.917 / 0.924 \\
jitter 2$u$ & 0.852 / 0.843 & 0.829 / 0.867 & 0.841 / 0.844 & 0.843 / 0.851 & 0.845 / 0.861 & 0.831 / 0.850 \\
\bottomrule
\end{tabular}
}
\end{table}

\begin{table}[!htbp]\centering\scriptsize
\caption{AOI-graph sensitivities on the two halves of a program reading under the default transformations and under support-matched transformations (redrawn until the trial keeps at least its observed dwell share minus 0.02; in brackets, the share of draws for which this succeeded within 200 tries), and the sensitivities of a descriptor made only of the fixation and dwell shares on the support. In the first column: observed AUC of the AOI graph and of the support-only descriptor. Bold: interval excludes zero.}
\resizebox{\linewidth}{!}{\begin{tabular}{lllll}
\toprule
Problem & Contrast & AOI graph & AOI graph, support-matched & Support-only descriptor \\
\midrule
\texttt{em\_halves} (AUC 0.682; 0.773) & rotation & \textbf{+0.098 [+0.031, +0.166]} & +0.032 [$-$0.059, +0.124] (99\%) & +0.054 [$-$0.001, +0.110] \\
\texttt{em\_halves} (AUC 0.682; 0.773) & jitter 0.25$u$ & +0.036 [$-$0.029, +0.102] & +0.019 [$-$0.048, +0.086] (100\%) & \textbf{+0.048 [+0.018, +0.078]} \\
\texttt{em\_halves} (AUC 0.682; 0.773) & jitter 0.5$u$ & +0.023 [$-$0.025, +0.071] & +0.035 [$-$0.029, +0.099] (100\%) & \textbf{+0.112 [+0.065, +0.159]} \\
\texttt{em\_halves} (AUC 0.682; 0.773) & jitter 1$u$ & +0.045 [$-$0.032, +0.123] & +0.057 [$-$0.027, +0.141] (89\%) & \textbf{+0.119 [+0.059, +0.178]} \\
\texttt{em\_halves} (AUC 0.682; 0.773) & jitter 2$u$ & \textbf{+0.113 [+0.014, +0.212]} & \textbf{+0.121 [+0.035, +0.207]} (31\%) & \textbf{+0.123 [+0.054, +0.192]} \\
\bottomrule
\end{tabular}
}
\end{table}

\begin{table}[!htbp]\centering\scriptsize
\caption{Duration-marginal transformation of the classic descriptor with the leave-one-participant-out donor pool of the main analysis and with training-fold pools (each fold's trials draw only from that fold's training participants; mean and SD over $R = 6$ replicates; within-fold AUC, class-balanced folds).}
\begin{tabular}{lrrrrr}
\toprule
Problem & AUC obs. & AUC, LOO pool & AUC, training-fold pool (SD) & $\Cs$ LOO & $\Cs$ training-fold \\
\midrule
\texttt{gb\_tex\_vd1} & 0.969 & 0.935 & 0.935 (0.000) & +0.035 & +0.034 \\
\texttt{gb\_vd1\_blg} & 0.886 & 0.852 & 0.852 (0.001) & +0.033 & +0.034 \\
\texttt{gb\_vd1\_vd2} & 0.640 & 0.516 & 0.523 (0.009) & +0.124 & +0.117 \\
\texttt{gb\_tex\_s1s2} & 0.539 & 0.508 & 0.503 (0.020) & +0.031 & +0.036 \\
\texttt{gb\_vd1\_halves} & 0.564 & 0.556 & 0.553 (0.004) & +0.008 & +0.011 \\
\texttt{sb\_native} & 0.689 & 0.672 & 0.672 (0.002) & +0.017 & +0.017 \\
\texttt{sb\_compr} & 0.587 & 0.561 & 0.562 (0.005) & +0.026 & +0.024 \\
\texttt{em\_program} & 0.539 & 0.507 & 0.532 (0.010) & +0.032 & +0.007 \\
\texttt{em\_expertise} & 0.438 & 0.461 & 0.474 (0.038) & $-$0.023 & $-$0.037 \\
\texttt{em\_correct} & 0.501 & 0.447 & 0.454 (0.023) & +0.054 & +0.047 \\
\texttt{em\_correctrect} & 0.523 & 0.474 & 0.411 (0.050) & +0.048 & +0.112 \\
\texttt{em\_correctveh} & 0.500 & 0.475 & 0.454 (0.070) & +0.025 & +0.046 \\
\texttt{em\_halves} & 0.776 & 0.772 & 0.772 (0.005) & +0.004 & +0.004 \\
\bottomrule
\end{tabular}

\end{table}

\begin{table}[!htbp]\centering\scriptsize
\caption{Three permutation distributions for negative sensitivities (within-fold AUC, class-balanced folds; 1000 permutations under the schemes of Table 2 of the article; $R = 6$ transformed replicates). The observed AUC, the transformed AUC and their contrast are each compared with their own permutation distribution; for the contrast the lower-tail probability is given. The first four rows are the order cells of the classic descriptor discussed in Section 5.4 of the article; the last is the unexplained cell of Section 5.11.}
\begin{tabular}{lllrrrrrr}
\toprule
Problem & Descriptor & Transformation & Observed AUC & $p$ & Transformed AUC & $p$ & Contrast & $p$ (lower tail) \\
\midrule
\texttt{em\_program} & Classic & order & 0.539 & 0.095 & 0.788 & 0.001 & $-0.249$ & 0.001 \\
\texttt{gb\_vd1\_halves} & Classic & order & 0.564 & 0.001 & 0.653 & 0.001 & $-0.090$ & 0.001 \\
\texttt{gb\_tex\_vd1} & Classic & order & 0.969 & 0.001 & 0.991 & 0.001 & $-0.022$ & 0.142 \\
\texttt{gb\_vd1\_blg} & Classic & order & 0.886 & 0.001 & 1.000 & 0.001 & $-0.114$ & 0.001 \\
\texttt{em\_halves} & RQA $0.5u$ & jitter $1u$ & 0.613 & 0.001 & 0.692 & 0.001 & $-0.080$ & 0.027 \\
\bottomrule
\end{tabular}

\end{table}

\begin{table}[!htbp]\centering\scriptsize
\caption{Calibration at the largest analysed sample size: 322 participants $\times$ 4 trials with the label varying within participants ($R = 6$, $B = 60$; 40 datasets per row; targets from 60 further datasets with $R = 20$, or zero by exchangeability for the order contrast). Coverage of the target by the 95\% interval; exclusions of zero where the target is zero by argument; mean width; mean bootstrap SE over the empirical SD of the estimate; coverage of the nominal 80\% interval.}
\begin{tabular}{lllrrrrr}
\toprule
Design & Descriptor & Contrast & Coverage & Excl.\ 0 at zero target & Width & SE ratio & Coverage at 80\% \\
\midrule
no signal & classic & order & 40/40 & 0/40 & 0.111 & 1.65 & 0.975 \\
no signal & directed & order & 38/40 & 2/40 & 0.139 & 1.39 & 0.950 \\
no signal & directed & jitter 160 px & 39/40 & -- & 0.148 & 1.26 & 0.900 \\
no signal & directed & jitter 40 px & 39/40 & -- & 0.150 & 1.39 & 0.975 \\
no signal & directed & rotation & 40/40 & -- & 0.144 & 1.42 & 0.900 \\
no signal & rqa & order & 40/40 & 0/40 & 0.133 & 1.63 & 0.950 \\
no signal & rqa & jitter 160 px & 40/40 & -- & 0.159 & 1.70 & 1.000 \\
no signal & rqa & jitter 40 px & 40/40 & -- & 0.150 & 1.62 & 0.925 \\
signal elsewhere & directed & order & 40/40 & 0/40 & 0.024 & 1.38 & 0.900 \\
signal elsewhere & directed & jitter 40 px & 40/40 & -- & 0.075 & 1.97 & 1.000 \\
signal elsewhere & directed & rotation & 40/40 & -- & 0.045 & 1.44 & 0.975 \\
signal elsewhere & rqa & order & 40/40 & 0/40 & 0.017 & 1.57 & 1.000 \\
signal elsewhere & rqa & jitter 40 px & 40/40 & -- & 0.070 & 1.81 & 0.975 \\
\bottomrule
\end{tabular}

\end{table}

\begin{table}[!htbp]\centering\scriptsize
\caption{Gradient-boosted trees on the cells of Table 7 of the article at the screening settings and at the long settings (seed 42). Bold: interval excludes zero.}
\begin{tabular}{lllll}
\toprule
Problem & Descriptor & Transformation & Trees, $R = 6$, $B = 60$ & Trees, $R = 20$, $B = 300$ \\
\midrule
\texttt{gb\_tex\_vd1} & Dur.\ trend & duration marginal & \textbf{+0.371 [+0.318, +0.424]} & \textbf{+0.378 [+0.327, +0.428]} \\
\texttt{gb\_vd1\_blg} & Classic & order & \textbf{$-$0.118 [$-$0.148, $-$0.088]} & \textbf{$-$0.118 [$-$0.146, $-$0.090]} \\
\texttt{gb\_vd1\_blg} & Dur.\ trend & duration marginal & \textbf{+0.226 [+0.161, +0.291]} & \textbf{+0.231 [+0.168, +0.293]} \\
\texttt{gb\_vd1\_vd2} & Classic & duration marginal & \textbf{+0.130 [+0.073, +0.187]} & \textbf{+0.127 [+0.074, +0.180]} \\
\texttt{gb\_vd1\_vd2} & Dispersion & jitter $2u$ & \textbf{+0.102 [+0.038, +0.166]} & \textbf{+0.100 [+0.036, +0.164]} \\
\texttt{gb\_vd1\_vd2} & Dur.\ trend & duration coupling & \textbf{+0.176 [+0.133, +0.219]} & \textbf{+0.172 [+0.128, +0.215]} \\
\texttt{gb\_vd1\_vd2} & Dur.\ trend & duration marginal & \textbf{+0.272 [+0.216, +0.329]} & \textbf{+0.274 [+0.226, +0.322]} \\
\texttt{gb\_vd1\_vd2} & Dur.\ trend & order & \textbf{+0.178 [+0.133, +0.224]} & \textbf{+0.175 [+0.132, +0.218]} \\
\texttt{gb\_vd1\_vd2} & RQA $0.5u$ & jitter $2u$ & \textbf{+0.106 [+0.041, +0.170]} & \textbf{+0.103 [+0.037, +0.168]} \\
\texttt{gb\_vd1\_vd2} & RQA $1u$ & jitter $2u$ & \textbf{+0.125 [+0.054, +0.195]} & \textbf{+0.123 [+0.067, +0.178]} \\
\texttt{gb\_vd1\_halves} & Directed & jitter $2u$ & \textbf{+0.192 [+0.140, +0.243]} & \textbf{+0.199 [+0.145, +0.252]} \\
\texttt{gb\_vd1\_halves} & Dispersion & jitter $2u$ & \textbf{+0.123 [+0.097, +0.150]} & \textbf{+0.124 [+0.096, +0.153]} \\
\texttt{gb\_vd1\_halves} & Dispersion & rotation & \textbf{+0.177 [+0.139, +0.215]} & \textbf{+0.167 [+0.133, +0.202]} \\
\texttt{gb\_vd1\_halves} & RQA $0.5u$ & jitter $2u$ & \textbf{+0.157 [+0.101, +0.214]} & \textbf{+0.165 [+0.108, +0.222]} \\
\texttt{gb\_vd1\_halves} & RQA $1u$ & jitter $2u$ & \textbf{+0.212 [+0.172, +0.253]} & \textbf{+0.219 [+0.168, +0.270]} \\
\texttt{sb\_native} & RQA $0.5u$ & jitter $2u$ & $-$0.002 [$-$0.084, +0.081] & +0.008 [$-$0.074, +0.091] \\
\texttt{sb\_native} & RQA $1u$ & jitter $2u$ & \textbf{+0.100 [+0.009, +0.190]} & \textbf{+0.111 [+0.024, +0.199]} \\
\texttt{sb\_native} & RQA $1u$ & jitter $1u$ & +0.071 [$-$0.003, +0.145] & \textbf{+0.082 [+0.012, +0.151]} \\
\texttt{em\_program} & Dispersion & rotation & \textbf{+0.166 [+0.088, +0.244]} & \textbf{+0.164 [+0.074, +0.255]} \\
\texttt{em\_halves} & Classic & order & \textbf{+0.128 [+0.033, +0.222]} & \textbf{+0.148 [+0.062, +0.233]} \\
\texttt{em\_halves} & RQA $0.5u$ & jitter $1u$ & $-$0.089 [$-$0.189, +0.011] & $-$0.087 [$-$0.192, +0.018] \\
\bottomrule
\end{tabular}

\end{table}

\clearpage
\section*{Methods appendix}
\subsection*{A1. Feature definitions}
Let a trial have fixations $(x_i, y_i, d_i)$, $i = 1, \dots, N$, in sequence order, with $d_i$ the duration.

\emph{Count.} $N$.

\emph{Dispersion.} The population SDs of $x$ and of $y$, and the mean pairwise distance $\binom{N}{2}^{-1} \sum_{i<j} \lVert p_i - p_j \rVert$, $p_i = (x_i, y_i)$.

\emph{Classic.} Mean duration $N^{-1} \sum_i d_i$ and mean saccade amplitude $(N - 1)^{-1} \sum_{i<N} \lVert p_{i+1} - p_i \rVert$.

\emph{Duration trend.} The Pearson correlation of $d_i$ with $i$, and the mean duration; trials with fewer than five fixations are excluded.

\emph{Recurrence} \citep{anderson2013recurrence}. Fixations $i < j$ recur if $\lVert p_i - p_j \rVert < r$; $n_r$ is the number of recurrent pairs, and the recurrence plot is the upper triangle of the $N \times N$ recurrence matrix. Recurrence rate: $100 \cdot 2 n_r / (N(N-1))$. Determinism: $100$ times the share of recurrent pairs that lie on diagonal lines of length at least 2. Laminarity: $100$ times the share that lie on vertical lines of length at least 2. Centre of recurrence mass: $100 \cdot \overline{(j - i)} / (N - 1)$, the mean over recurrent pairs. A trial without recurrent pairs has all four equal to zero; trials with fewer than five fixations are excluded. Radii: $r = 0.5u$ and $1u$.

\emph{Directed transition graph.} Each fixation is assigned to the tile $(\lfloor x/u \rfloor, \lfloor y/u \rfloor)$; consecutive repeats are collapsed to a sequence $s$ over the visited tiles $V$; $A_{ab}$ counts immediate transitions $a \to b$, and $E = \{(a, b): a \neq b, A_{ab} > 0\}$. Reciprocity: the share of edges in $E$ whose reverse is in $E$. Revisit tendency: $1 - |V|/|s|$. Transition entropy: $-\sum_{ab} P_{ab} \log_2 P_{ab}$ with $P = A / \sum A$. Degree asymmetry: the mean over tiles of $|\mathrm{out}_a - \mathrm{in}_a| / (\mathrm{out}_a + \mathrm{in}_a)$, with out- and in-strengths the row and column sums of $A$.

\emph{Fixation Adjacency Graph.} Anchors are triangulated (Delaunay). The dwell of a triangle is the summed duration of the fixations inside it; fixations in no triangle are dropped. Vertices are the triangles with positive dwell; two are adjacent if they share an anchor; the weight of an edge is the summed dwell of its two triangles. Attention centrality: the largest, over connected components, of the weight of the component's maximum spanning tree, divided by the total edge weight (set to 0.25 for a graph without edges). Betweenness: the maximum edge betweenness centrality \citep{brandes2001faster}. Compactness: the Estrada index $\sum_k e^{\lambda_k}$ of the adjacency matrix \citep{estrada2000characterization}, averaged over components. Squareness: for each component, $\Delta_x \Delta_y / \max(\Delta_x, \Delta_y)^2$ with $\Delta_x, \Delta_y$ the extents of its triangle centroids (zero for a single triangle), averaged over components. The fifth feature is the number of components. Lattice anchors are $2u$ apart and cover the disc swept by the display under rotation; AOI anchors are the token centres of the program; the padded AOI graph adds lattice anchors $2u$ apart around them, leaving out those closer than $2u$ to a token anchor.

\subsection*{A2. Class sizes and fixation counts}
\begin{table}[H]\centering\scriptsize
\caption{Trials and participants per class, and fixations per trial, for every problem (after the exclusions of Section 4 of the article). Class labels are those of the released loaders.}
\begin{tabular}{llrrrrrr}
\toprule
 & & \multicolumn{2}{c}{Trials} & \multicolumn{2}{c}{Participants} & \multicolumn{2}{c}{Fixations per trial} \\
\cmidrule(lr){3-4}\cmidrule(lr){5-6}\cmidrule(lr){7-8}
Problem & Positive class & positive & negative & positive & negative & median (IQR) & range \\
\midrule
\texttt{gb\_tex\_vd1} & \texttt{VD1} & 644 & 644 & 322 & 322 & 179 (132--225) & 37--300 \\
\texttt{gb\_vd1\_blg} & \texttt{BLG} & 643 & 644 & 322 & 322 & 121 (96--146) & 33--392 \\
\texttt{gb\_vd1\_vd2} & \texttt{VD2} & 644 & 644 & 322 & 322 & 141 (121--159) & 32--233 \\
\texttt{gb\_vd1\_halves} & \texttt{VD1\#2} & 644 & 644 & 322 & 322 & 66 (57--75) & 18--113 \\
\texttt{gb\_tex\_s1s2} & \texttt{TEX@S2} & 322 & 322 & 322 & 322 & 225 (206--243) & 97--300 \\
\texttt{sb\_native} & \texttt{nonnative} & 620 & 1426 & 29 & 66 & 122 (92--152) & 7--398 \\
\texttt{sb\_compr} & \texttt{low} & 805 & 558 & 79 & 66 & 121 (93--152) & 8--398 \\
\texttt{em\_program} & \texttt{vehicle} & 112 & 112 & 112 & 112 & 218 (145--291) & 25--868 \\
\texttt{em\_halves} & \texttt{second} & 224 & 224 & 112 & 112 & 109 (72--145) & 12--434 \\
\texttt{em\_expertise} & \texttt{low} & 64 & 51 & 37 & 28 & 181 (129--263) & 25--634 \\
\texttt{em\_correct} & \texttt{incorrect} & 147 & 124 & 120 & 108 & 212 (142--286) & 25--868 \\
\texttt{em\_correctrect} & \texttt{incorrect} & 41 & 92 & 41 & 92 & 193 (142--279) & 35--606 \\
\texttt{em\_correctveh} & \texttt{incorrect} & 106 & 32 & 106 & 32 & 230 (143--293) & 25--868 \\
\bottomrule
\end{tabular}

\end{table}

\subsection*{A3. Software and settings}
Python 3.13.16; NumPy 2.5.3; SciPy 1.18.1; pandas 3.0.5; scikit-learn 1.7.2; XGBoost 3.2.0; NetworkX 3.7; package \texttt{scanpathdecomp} 0.6.3. Linear SVM: \texttt{LinearSVC} (liblinear; squared hinge loss, L2 penalty, $C = 1$, at most 20\,000 iterations, random state 0) on features standardized with training-fold means and SDs. Gradient-boosted trees: \texttt{XGBClassifier} with 200 trees of depth 3, learning rate 0.05, row subsampling 0.8, histogram tree method, logistic loss, random state 0. Folds: ten participant folds, participants dealt within class in a seeded order (seed 0 for the reported profiles; seeds 0--19 for the partition check).
\bibliographystyle{apalike}
\bibliography{references}
\end{document}
